\documentclass[11pt,a4paper]{article}
\usepackage[margin=25mm,headheight=14pt]{geometry}
\usepackage{fontspec}
\setmainfont{Latin Modern Roman}
\setsansfont{Latin Modern Sans}
\setmonofont{Latin Modern Mono}[Scale=0.86]
\usepackage{amsmath,amssymb}
\usepackage{unicode-math}
\setmathfont{Latin Modern Math}
\usepackage{microtype}
\usepackage{graphicx}
\usepackage{longtable,booktabs,array}
\newcounter{none}
\usepackage{calc}
\usepackage{xcolor}
\definecolor{linkblue}{HTML}{244B70}
\usepackage{fvextra}
\DefineVerbatimEnvironment{verbatim}{Verbatim}{fontsize=\small,breaklines=true,breakanywhere=true,breaksymbolleft={},breaksymbolright={},xleftmargin=5pt}
\usepackage{enumitem}
\setlist{itemsep=2pt,parsep=1pt,topsep=4pt}
\usepackage{fancyhdr}
\pagestyle{fancy}
\fancyhf{}
\fancyhead[L]{\small SU(2) lattice Yang-Mills heat correlations}
\fancyhead[R]{\small 21 September 2026}
\fancyfoot[C]{\thepage}
\renewcommand{\headrulewidth}{0.3pt}
\usepackage{hyperref}
\hypersetup{colorlinks=true,linkcolor=linkblue,urlcolor=linkblue,pdftitle={Heat
correlations and the full physical complement in SU(2) lattice
Yang-Mills
theory},pdfauthor={KokunoYumeto},pdfsubject={Fourth-order heat, volume-independent error, and the full physical complement},pdfcreator={Pandoc and XeLaTeX}}
\usepackage{bookmark}
\urlstyle{same}
\setlength{\emergencystretch}{4em}
\hfuzz=0.2pt
\setlength{\parindent}{0pt}
\setlength{\parskip}{5pt plus 1pt minus 1pt}
\setcounter{secnumdepth}{2}
\setcounter{tocdepth}{2}
\providecommand{\tightlist}{\setlength{\itemsep}{0pt}\setlength{\parskip}{0pt}}
\providecommand{\pandocbounded}[1]{#1}
\let\originalsection\section
\renewcommand{\section}{\clearpage\originalsection}
\let\originaldisplaystart\[
\let\originaldisplayend\]
\renewcommand{\[}{\begingroup\small\originaldisplaystart}
\renewcommand{\]}{\originaldisplayend\endgroup}
\begin{document}
\begin{titlepage}
\vspace*{14mm}
{\sffamily\Large MATHEMATICAL READER\par}
\vspace{12mm}
{\LARGE\bfseries Heat correlations and the full physical complement in
SU(2) lattice Yang-Mills theory\par}
\vspace{9mm}
{\large Fourth-order coefficients, volume-independent bounds, and
original state and energy metrics\par}
\vspace{13mm}
{\large KokunoYumeto\par}
\vspace{3mm}
{21 September 2026\par}
\vfill
\begin{minipage}{0.96\textwidth}
A mathematical introduction followed by all five complete proof manuscripts. The original link coordinates, Haar pairings, physical time and complementary states are retained. The results concern a stated strong-coupling domain and spatial volume at fixed lattice spacing; a four-dimensional continuum mass gap is not claimed.

\medskip
The earlier workbench's distinct \(S^6\) and Jacobian inputs are credited to Levent Alpöge and their named collaborators in the introduction. The present heat transfer's precise dependencies are identified separately.
\end{minipage}
\vspace{10mm}
\end{titlepage}
\pagenumbering{roman}
\tableofcontents
\clearpage
\pagenumbering{arabic}
\section{Reading the heat and volume
results}\label{reading-the-heat-and-volume-results}

The basic question of this edition is concrete. A plaquette observable
probes a small square of a lattice gauge field. How does its correlation
with another plaquette decay under the interacting heat flow? Can one
control the error in a fourth-order calculation without that error
growing with the size of the lattice? Finally, when states built from
those observables are allowed to mix with every other physical state,
how much can their energy change?

The five manuscripts collected here answer these questions on an
explicitly specified strong-coupling domain. They give the complete
fourth-order heat coefficients, a bound on their all-order error that is
independent of spatial volume, and quantitative control of the entire
complementary physical space. The volume-independent estimates also
construct an infinite spatial lattice at \textbf{fixed lattice spacing}.
They do not establish a four-dimensional continuum Yang-Mills field or
its continuum mass gap.

This introduction states the objects and results and explains how the
proofs fit together. Appendices A-E reproduce all five supplied proof
texts, including their original equation labels, parameter restrictions,
and source citations. The exact coefficient tables, proofs, replay
scripts and historical records are in the
\href{https://github.com/KokunoYumeto/yang-mills-interacting-workbench/tree/main/yang-mills/consolidation/20260921}{public
workbench}. This is a readable companion to those records, not a
replacement for the complete calculation.

\subsection{The original lattice and its physical Hilbert
space}\label{the-original-lattice-and-its-physical-hilbert-space}

Let \(L\geq2\) be an integer. The vertices are \(\{-L,\ldots,L\}^3\);
all contained positive nearest-neighbour links and elementary plaquettes
are retained. A link carries a matrix \(U_e\in\mathrm{SU}(2)\). The
Hilbert space is the subspace of the product-Haar \(L^2\) space
invariant under gauge transformations at every vertex, including the
boundary vertices. Its inner product is conjugate-linear in the first
entry.

For \(p=(n;i,j)\), \(i<j\), the plaquette observable is the fundamental
trace

\[
W_p=\operatorname{tr}\!\left(U_i(n)U_j(n+e_i)
U_i(n+e_j)^{-1}U_j(n)^{-1}\right).
\]

With \(T_\alpha=-i\sigma_\alpha/2\) and the corresponding original link
derivatives \(X_{e,\alpha}\), set

\[
K_L=-\sum_{e,\alpha}X_{e,\alpha}^2,\qquad
H_L=\kappa K_L+\kappa\xi\sum_{p\in\mathsf P_L}(2-W_p),
\qquad \kappa=\frac{2g^2}{a},\quad \xi=\frac1{4g^4}.
\]

Here \(a>0\) is lattice spacing and \(g>0\) is coupling. The operator
and form domains are the physical parts of \(H^2\) and \(H^1\) on the
finite product of link groups. The smooth positive unit ground vector is
\(\psi_L\), its energy is \(E_{0,L}\), and

\[
A_L=H_L-E_{0,L},\qquad \rho_L=\psi_L^2,\qquad
r_p=(W_p-\langle W_p\rangle_{\rho_L})\psi_L.
\]

Thus \(r_p\) is centered in the \emph{interacting} vacuum, not just in
Haar measure. On the centered physical space
\(\mathcal H_{0,L}=\psi_L^\perp\), the connected heat correlation is

\[
\widehat C_{L,pq}(\tau;\xi)
=\langle r_p,e^{-\tau A_L/\kappa}r_q\rangle,
\qquad \tau=\kappa t.
\]

The physical time is \(t\); \(\tau\) is its stated dimensionless
coordinate. The number of plaquettes is \(M_L=3(2L)^2(2L+1)\), so
\(M_2=240\).

\subsection{What the fourth-order calculation
contains}\label{what-the-fourth-order-calculation-contains}

The link-center symmetry makes the two-mark correlation even in the
homogeneous source \(\xi\). The calculation therefore has the form

\[
\widehat C_L(\tau;\xi)
=C_{0,L}(\tau)+\xi^2C_{2,L}(\tau)+\xi^4C_{4,L}(\tau)
+\mathcal R_{6,L}(\tau;\xi).
\]

Appendix A, H2-H4, computes the coefficients from the actual
ground-state and resolvent recurrences. The ground-energy shift,
division by the original vacuum norm, and subtraction of the ground pole
are included. Each connected Laplace coefficient is a finite
partial-fraction sum

\[
F_\nu(z)=\sum_{c>0,\,m\geq1}\frac{a_{c,m}}{(z+c)^m},
\qquad
C_\nu(\tau)=\sum_{c,m}
\frac{a_{c,m}}{(m-1)!}\tau^{m-1}e^{-c\tau}.
\]

There are 84 exact time functions across 17 geometric cases. Their
\textbf{Laplace transforms} are rational functions; the time functions
are finite polynomial-exponential sums. The phrase "rational time
functions" in the preserved source should be read with that distinction.
On the \(L=2\) box the resulting coefficient matrices have respectively
240, 2,496 and 9,660 nonzero entries at degrees zero, two and four.
Their time integrals return the earlier static response matrices.

The finite first excitation band is a related, but separate, result. If
\(Y_\xi\) is the actual projected plaquette frame, its state Gram
\(G_{\rm band}=Y_\xi^*Y_\xi\) need not be the identity. Appendix A, H8,
gives

\[
A_{\rm band}=3I+\xi^2T_2+\xi^4T_4+\cdots,
\qquad T_2=\frac7{15}I-\frac{D_{\rm degree}+A_{\rm adj}}{21},
\]

\[
T_4-T_4^*=T_2G_{{\rm band},2}-G_{{\rm band},2}T_2.
\]

This last identity explains why an ordinary coordinate adjoint is not
the adjoint in the computed band metric. On the open \(L=2\) box its
right-hand side has 1,440 nonzero entries. The band remainder is proved
only for \(|\xi|<1/(32M_L)\); the later volume-independent \emph{heat}
result does not silently enlarge that band domain.

\subsection{From a finite-box estimate to a volume-independent heat
bound}\label{from-a-finite-box-estimate-to-a-volume-independent-heat-bound}

Appendix B first proves an entrywise analytic bound on \(|\xi|<1/M_L\).
This controls a fixed box, but its source radius shrinks as the box
grows. Appendix D addresses exactly that dependence. The proof uses a
locally labelled Fourier-coefficient norm to construct the exponential
source, while keeping the original Haar pairing and all scalar energy
terms.

The essential estimates are as follows. The physical representation
labels satisfy \(c_j\geq6j_e\) at each active link and \(c_j\geq3\) away
from the constant block. The original plaquette coefficient has trace
norm \(8\). These facts give a convergent labelled source series on

\[
R_0=\frac3{256},\qquad
\chi(r)=\frac{1-\sqrt{1-256r/3}}2,\qquad
d(r)=3(1-\chi(r)).
\]

The scalar mean is solved together with the nonconstant equation. A
closed generator on a stated dense coefficient domain then gives the
actual heat flow, rather than replacing a semigroup estimate by an
estimate on one inverse. The Poisson identity pairs that flow against an
arbitrary signed sum of plaquettes. Taking the supremum over the signs
controls an entire matrix row.

For \(\|B\|_{\rm row}=\max_p\sum_q|B_{pq}|\), the result is

\[
\boxed{\displaystyle
\|\mathcal R_{6,L}(\tau;\xi)\|_{\rm row}
\leq512e^{-3\tau/2}\frac{\theta^6}{1-\theta^2},
\qquad \theta=\frac{|\xi|}{R_0}<1.}
\]

It holds for every \(L\geq2\) and \(\tau\geq0\) with the same constants
(Appendix D, U29-U30). A second bound, U29b, has the smaller radius
\(R_1=45/4096\), prefactor \(6184/25\), and decay \(15/8\). Both bounds
are retained with their own radii; the better of their scalar error
bounds may be used where both apply.

Local Taylor coefficients agree once their original support fits inside
both boxes. Combined with the complete analytic tail, this gives
full-sequence convergence of local vacuum means and heat correlations.
U9 constructs the limiting gauge-invariant state and symmetric Markov
semigroup, with a self-adjoint generator on its physical \(L^2\) space.
This passage increases spatial volume with \(a\) fixed. It is not an
ultraviolet limit \(a\to0\).

\subsection{The three Grams and why the complement
matters}\label{the-three-grams-and-why-the-complement-matters}

Let \(R:\mathbb C^{M_L}\to\mathcal H_{0,L}\) be the column map
\(Rc=\sum_pc_pr_p\), and define its inverse-energy family

\[
\Phi=\kappa A_L^{-1}R,\qquad A_L\Phi=\kappa R.
\]

The three original Grams measure different things:

\[
G_0=R^*R,\qquad G_1=R^*\Phi,\qquad
G_2=\Phi^*\Phi,\qquad E=\Phi^*A_L\Phi=\kappa G_1.
\]

Here \(G_0\) is the norm of the forcing columns, \(G_2\) is the norm of
the inverse-energy states, and \(E\) is their energy form. More
generally, \(G_k=\kappa^kR^*A_L^{-k}R\) for \(k\geq1\). The heat
integral computes all these inverse-energy moments:

\[
G_k=\frac1{(k-1)!}\int_0^\infty
\tau^{k-1}\widehat C_L(\tau;\xi)\,d\tau\quad(k\geq1).
\]

At \(g^2\geq16\), Appendix E, M10, obtains, uniformly over \(L\geq2\)
and \(a>0\),

\[
\|G_0-I\|<\frac{124}{10^6},\qquad
\|G_1-I/3\|<\frac{66}{10^6},\qquad
\|G_2-I/9\|<\frac{36}{10^6}.
\]

The coefficient bound is valid at boundary rows because absolute values
are summed \emph{before} contributions from distinct supports can
cancel. The full-space estimate is \(A_L\geq(117/40)\kappa I\) on
\(\mathcal H_{0,L}\).

These facts alone would not justify pretending that
\(\operatorname{im}\Phi\) is a reducing subspace. Its coupling to the
rest of the Hilbert space must be retained. Define

\[
P=\Phi G_2^{-1}\Phi^*,\qquad Q=I-P,\qquad B=QR.
\]

Then \(P\) is the actual orthogonal state projection and

\[
B^*B=G_0-G_1G_2^{-1}G_1.
\]

For every form-domain vector written as \(\Phi c+h\),
\(h\in Q\mathcal H_{0,L}\), the full energy is

\[
q_{A_L}(\Phi c+h)
=c^*Ec+2\kappa\operatorname{Re}\langle Bc,h\rangle+q_{A_L}(h).
\]

Appendix E, M5, proves that the complementary operator \(D_Q=QA_LQ\) on
\(\operatorname{Dom}(A_L)\cap Q\mathcal H_{0,L}\) is self-adjoint and
bounded below. Minimization over its entire form domain therefore gives

\[
h_{\min}(c)=-\kappa D_Q^{-1}Bc,\qquad
E_{\rm full}=E-\kappa^2B^*D_Q^{-1}B,
\]

\[
G_{\rm restored}=G_2+\kappa^2B^*D_Q^{-2}B.
\]

At \(g^2\geq16\) the sharpened bounds are

\[
\frac{999}{1000}E\prec E_{\rm full}\preceq E,
\qquad
G_2\preceq G_{\rm restored}\prec\frac{1001}{1000}G_2.
\]

Thus the original family\textquotesingle s energy decreases by less than
one part in a thousand when the entire complement is allowed to relax,
while its state Gram increases by less than one part in a thousand. The
inequalities are between quadratic forms in the same original
coefficients; a strict inequality is evaluated on nonzero coefficients.
They are not quotients of matrix entries. M8 carries the bounded maps
and estimates to \(\ell^2\) of the infinite plaquette family, including
its full physical complement.

\subsection{Two concrete consequences}\label{two-concrete-consequences}

First, take opposite faces \(p=(0,0,0;0,1)\) and \(q=(0,0,1;0,1)\).
Their degree-zero and degree-two correlations vanish; the degree-four
term consists of four original three-face paths and one cube. At
\(\xi=10^{-10}\), so that \(g^2=50000\) and \(\kappa=100000/a\),
Appendix E, M26, gives

\[
\frac{85910}{10^7}
<\frac{C_{pq}(1/\kappa;\xi)}{\xi^4}
<\frac{85920}{10^7}.
\]

The complete correlation is positive throughout
\(1/\kappa\leq t\leq3/\kappa\), in every \(L\geq2\) and in the
fixed-spacing spatial limit. The all-order remainder is smaller than the
explicitly proved coefficient lower bound on this interval. Positivity
is therefore a statement about the actual interacting correlation, not
merely its truncated expansion.

Second, finite heat time gives a controlled replacement of
inverse-energy states:

\[
\Phi_T=(I-e^{-TA_L})\Phi
=\kappa\int_0^Te^{-tA_L}R\,dt,\qquad
A_L\Phi_T=\kappa R-\kappa e^{-TA_L}R.
\]

The omitted forcing is explicit. If \(\varepsilon=e^{-\kappa dT}\) with
the full-space lower bound \(A_L\geq\kappa dI\), spectral calculus gives

\[
(1-\varepsilon)^2G_2\preceq\Phi_T^*\Phi_T\preceq G_2,
\qquad
(1-\varepsilon)^2E\preceq\Phi_T^*A_L\Phi_T\preceq E.
\]

These inequalities survive restriction and minimization over the
\emph{same coefficient fibers}. In particular, a signed sum
\(\mathcal L=\sum_ja_j\log\det Q_j\) of such original positive quotient
forms satisfies

\[
|\mathcal L_T-\mathcal L|
\leq2\mathcal B[-\log(1-\varepsilon)],\qquad
\mathcal B=\sum_j|a_j|\operatorname{rank}Q_j.
\]

For four returns with coefficients \(a_j\in\{-1,+1\}\) and ranks at most
240, so that \(\mathcal B\leq960\), \(g^2\geq16\) and \(T=10/\kappa\)
give error below \(4\times10^{-10}\) in every exterior volume. For a
growing observation family, M25 keeps the rank explicitly and reaches
tolerance \(\eta>0\) at

\[
T=\frac1{\kappa d}\log\!\left(1+\frac{2\mathcal B}{\eta}\right).
\]

\subsection{Sources, mathematical lineage, and how to verify the
edition}\label{sources-mathematical-lineage-and-how-to-verify-the-edition}

The exponential-vacuum and character framework has a human antecedent in
Schütte, Zheng Weihong and Hamer {[}1{]}. The block minimization is a
Feshbach-Schur construction; {[}2{]} is a modern primary reference.
Appendix D proves the needed compact-group Fourier coefficient estimates
and cites Eymard {[}3{]}. Its generator argument states the domain,
dissipativity and range conditions used from Lumer and Phillips {[}4{]}.
These are point-of-use dependencies, not generic acknowledgments.

Two distinct contributions associated with \textbf{Levent Alpöge} enter
the earlier workbench and deserve separate attribution. The \(S^6\)
branch starts from the claimed complex-threefold construction he
circulated with AI assistance {[}5{]}, also credited to him in
Engel\textquotesingle s exposition {[}6{]}. A different branch starts
from his Jacobian counterexample: global noninjectivity despite a
constant nonzero Jacobian. His original announcement {[}7{]} credits
Akhil\textquotesingle s question and Fable\textquotesingle s work;
Tao\textquotesingle s later exposition is credited separately {[}8{]}.
These constructions are not interchangeable. Their subsequent coordinate
maps and conductor calculations belong to their respective
continuations. The five heat manuscripts collected here do not directly
use either originating construction as a lemma: their actual transferred
inputs are the heat, Gram, quotient-minimum and inverse-power maps
specified next. This distinction records the programme\textquotesingle s
mathematical origins without attributing the present heat bounds to
Alpöge or claiming an unproved spectral equivalence.

The workbench\textquotesingle s
\href{https://github.com/KokunoYumeto/yang-mills-interacting-workbench/blob/main/ATTRIBUTION.md}{attribution
and mathematical lineage record} gives the primary sources and exact
historical points of use.

The direct transfer used in the present proofs is the heat/Gram and
minimum-fiber argument from the Riemann workbench, not a spectral
identification of its arithmetic objects with Yang-Mills fields.
Appendix C identifies the receiving maps explicitly. Its pinned source
revision is
{\ttfamily 1\allowbreak{}2\allowbreak{}8\allowbreak{}a\allowbreak{}a\allowbreak{}3\allowbreak{}0\allowbreak{}8\allowbreak{}d\allowbreak{}d\allowbreak{}2\allowbreak{}a\allowbreak{}7\allowbreak{}0\allowbreak{}b\allowbreak{}a\allowbreak{}6\allowbreak{}e\allowbreak{}0\allowbreak{}7\allowbreak{}3\allowbreak{}8\allowbreak{}1\allowbreak{}6\allowbreak{}a\allowbreak{}9\allowbreak{}5\allowbreak{}7\allowbreak{}d\allowbreak{}1\allowbreak{}e\allowbreak{}e\allowbreak{}6\allowbreak{}4\allowbreak{}1\allowbreak{}4\allowbreak{}b\allowbreak{}a\allowbreak{}2\allowbreak{}c\allowbreak{}}
in
\href{https://github.com/KokunoYumeto/zeta-function-research-reader/tree/128aa308dd2a70ba6e073816a957d1ee6414ba2c}{the
Riemann workbench}: HM6-HM14 and HM19-HM24 of
{\ttfamily H\allowbreak{}E\allowbreak{}A\allowbreak{}T\allowbreak{}\_\allowbreak{}T\allowbreak{}O\allowbreak{}\_\allowbreak{}O\allowbreak{}R\allowbreak{}I\allowbreak{}G\allowbreak{}I\allowbreak{}N\allowbreak{}A\allowbreak{}L\allowbreak{}\_\allowbreak{}Q\allowbreak{}U\allowbreak{}O\allowbreak{}T\allowbreak{}I\allowbreak{}E\allowbreak{}N\allowbreak{}T\allowbreak{}\_\allowbreak{}M\allowbreak{}E\allowbreak{}T\allowbreak{}R\allowbreak{}I\allowbreak{}C\allowbreak{}S\allowbreak{}.\allowbreak{}t\allowbreak{}e\allowbreak{}x\allowbreak{}},
HG6-HG10 of
{\ttfamily O\allowbreak{}R\allowbreak{}I\allowbreak{}G\allowbreak{}I\allowbreak{}N\allowbreak{}A\allowbreak{}L\allowbreak{}\_\allowbreak{}H\allowbreak{}E\allowbreak{}A\allowbreak{}T\allowbreak{}\_\allowbreak{}A\allowbreak{}N\allowbreak{}D\allowbreak{}\_\allowbreak{}S\allowbreak{}O\allowbreak{}U\allowbreak{}R\allowbreak{}C\allowbreak{}E\allowbreak{}\_\allowbreak{}E\allowbreak{}N\allowbreak{}D\allowbreak{}P\allowbreak{}O\allowbreak{}I\allowbreak{}N\allowbreak{}T\allowbreak{}.\allowbreak{}t\allowbreak{}e\allowbreak{}x\allowbreak{}},
and IK9-IK12 of
{\ttfamily I\allowbreak{}N\allowbreak{}V\allowbreak{}E\allowbreak{}R\allowbreak{}S\allowbreak{}E\allowbreak{}\_\allowbreak{}P\allowbreak{}O\allowbreak{}W\allowbreak{}E\allowbreak{}R\allowbreak{}\_\allowbreak{}C\allowbreak{}O\allowbreak{}N\allowbreak{}D\allowbreak{}U\allowbreak{}C\allowbreak{}T\allowbreak{}O\allowbreak{}R\allowbreak{}\_\allowbreak{}O\allowbreak{}B\allowbreak{}S\allowbreak{}E\allowbreak{}R\allowbreak{}V\allowbreak{}A\allowbreak{}B\allowbreak{}I\allowbreak{}L\allowbreak{}I\allowbreak{}T\allowbreak{}Y\allowbreak{}.\allowbreak{}t\allowbreak{}e\allowbreak{}x\allowbreak{}}.
The source intake records retain their full paths and blob identities.
Original conductor coefficients, zeta-zero data and arithmetic constants
remain attached to their source constructions.

The workbench was developed with AI assistance from ChatGPT 5.6 Sol in
Ultra mode in Codex and GPT-6 Astra; the collection preserves the
individual mathematical and source records. The supplied mathematical
arguments are presented for scrutiny. Fresh integration checks repeated
2,425 exact volume checks with 16 false-formula controls, 879 heat
checks with 21 controls, and 52 route calculations, in both ordinary and
optimized Python. The cumulative manifest passed before and after these
runs. These executions verify their declared finite algebra and
endpoints. They are not Lean proofs, independent external peer review,
or a formal verification of every infinite-dimensional analytical step.

The source texts remain complete below, so the numerical statements can
be followed through their constructions rather than accepted from this
introduction. The original cumulative volume archive contains 2,015
files and retains all files of the preceding heat edition, including the
five earlier control documents in its history directory. A non-mutating
integrity check from the extracted archive root is
{\ttfamily p\allowbreak{}y\allowbreak{}t\allowbreak{}h\allowbreak{}o\allowbreak{}n\allowbreak{} \allowbreak{}-\allowbreak{}B\allowbreak{} \allowbreak{}v\allowbreak{}e\allowbreak{}r\allowbreak{}i\allowbreak{}f\allowbreak{}y\allowbreak{}\_\allowbreak{}c\allowbreak{}u\allowbreak{}m\allowbreak{}u\allowbreak{}l\allowbreak{}a\allowbreak{}t\allowbreak{}i\allowbreak{}v\allowbreak{}e\allowbreak{}.\allowbreak{}p\allowbreak{}y\allowbreak{}}.
Full producer replays should use a new disposable copy; their scripts
generate files and are not read-only checks.

The mathematical advance described here is the passage from finite-box
heat coefficients to a volume-independent heat error and to a
quantitative full-complement return in original physical metrics. The
next, different issue is extension beyond the stated strong-coupling
domain. The simultaneous shrinking-spacing path discussed in the source
eventually leaves that domain. Neither the existing local source
certificate nor its elementary restart establishes the ultraviolet
continuum limit.

\subsection{References}\label{references}

\begin{enumerate}
\def\labelenumi{\arabic{enumi}.}
\item
  D. Schütte, Zheng Weihong and C. J. Hamer, \emph{The Coupled Cluster
  Method in Hamiltonian Lattice Field Theory}, 1996,
  \href{https://arxiv.org/abs/hep-lat/9603026}{{a\allowbreak{}r\allowbreak{}X\allowbreak{}i\allowbreak{}v\allowbreak{}:\allowbreak{}h\allowbreak{}e\allowbreak{}p\allowbreak{}-\allowbreak{}l\allowbreak{}a\allowbreak{}t\allowbreak{}/\allowbreak{}9\allowbreak{}6\allowbreak{}0\allowbreak{}3\allowbreak{}0\allowbreak{}2\allowbreak{}6\allowbreak{}}}.
\item
  G. Dusson, I. M. Sigal and B. Stamm, \emph{The Feshbach-Schur map and
  perturbation theory}, 2021,
  \href{https://arxiv.org/abs/2105.02058}{arXiv:2105.02058}.
\item
  P. Eymard, \emph{L\textquotesingle algèbre de Fourier
  d\textquotesingle un groupe localement compact}, Bulletin de la
  Société Mathématique de France 92 (1964), 181-236.
  \href{https://www.numdam.org/item/BSMF_1964__92__181_0/}{Original
  article}.
\item
  G. Lumer and R. S. Phillips, \emph{Dissipative operators in a Banach
  space}, Pacific Journal of Mathematics 11 (1961), 679-698, Theorem
  3.1, pp. 686-687.
  \href{https://msp.org/pjm/1961/11-2/pjm-v11-n2-p19-s.pdf}{Original
  article}.
\item
  Levent Alpöge, \(S^6\) construction circulated with AI assistance,
  \href{https://alpo.ge/s6.pdf}{original 108-page manuscript},
  especially p. 2, Setup, for the coordinate matrix \(\Pi\).
\item
  Philip Engel, \emph{Complex structures on \(S^6\)}, 13 September 2026,
  \href{https://philip-engel.github.io/S6.pdf}{expository manuscript},
  abstract and Section 1.3.
\item
  Levent Alpöge,
  \href{https://x.com/__alpoge__/status/2079028340955197566}{original
  announcement of the Jacobian counterexample}, 20 July 2026. The
  announcement credits Akhil and Fable.
\item
  Terence Tao, \emph{A digestion of the Jacobian conjecture
  counterexample}, 21 July 2026,
  \href{https://terrytao.wordpress.com/2026/07/21/a-digestion-of-the-jacobian-conjecture-counterexample/}{exposition}.
\end{enumerate}

\clearpage
\appendix

\section{Original Yang--Mills heat correlations: the complete
fourth-degree
return}\label{original-yangmills-heat-correlations-the-complete-fourth-degree-return}

21 September 2026. This is an additive continuation of the recovered
cumulative workbench. The original open-box SU(2) operator, physical
units, Haar measure, gauge constraints and source coordinates remain.
The coefficient calculation uses the complete available fifth-source
archive. The separately saved sixth-source proofs are preserved with
their recovery status; their missing catalogues are not prerequisites
for this calculation. General exponential-vacuum and character methods
retain Schütte--Zheng--Hamer (1996) as a human antecedent. The exact
finite calculations here are accompanied by analytic proofs and
executable audits; independent external analytical review and historical
priority are not asserted.

\subsection{H1. The actual operator and time
coordinate}\label{h1.-the-actual-operator-and-time-coordinate}

For L\textgreater=2 use every vertex n in \{-L,...,L\}\^{}3, every
contained positive edge (n,i), and every contained elementary face
p=(n;i,j), i\textless j. Put

\begin{verbatim}
W_p=tr(U_i(n) U_j(n+e_i) U_i(n+e_j)^(-1) U_j(n)^(-1)),
T_alpha=-i sigma_alpha/2,
X_e,alpha f=d/ds f(...,exp(s T_alpha)U_e,...) at s=0,
K=-sum_(e,alpha) X_e,alpha^2,
H(x)=kappa K+kappa sum_p x_p(2-W_p),
kappa=2g^2/a,   x_p=xi=1/(4g^4) on the homogeneous line.       (H1)
\end{verbatim}

The inner product on the original product Haar probability space is
conjugate linear in its first entry. The physical subspace is invariant
under all vertex gauge transformations, including boundary vertices. The
domains are the physical parts of H\^{}2 and H\^{}1. The original
bounded smooth potential preserves these domains; the compact elliptic
operator has a smooth strictly positive unit ground state psi\_x with
energy E0(x). These domain and vacuum facts are the ones in the pinned
original finite-box source
{(\allowbreak{}S\allowbreak{}O\allowbreak{}U\allowbreak{}R\allowbreak{}C\allowbreak{}E\allowbreak{}\_\allowbreak{}I\allowbreak{}N\allowbreak{}T\allowbreak{}A\allowbreak{}K\allowbreak{}E\allowbreak{}.\allowbreak{}j\allowbreak{}s\allowbreak{}o\allowbreak{}n\allowbreak{},\allowbreak{}}
YM0).

The coefficient operator and its inverse physical return are exactly

\begin{verbatim}
B_x=(H(x)-2kappa sum_p x_p I)/kappa=K-sum_p x_p W_p,
H(x)=kappa B_x+2kappa sum_p x_p I,
E0(x)=kappa[2sum_p x_p+e(x)].                              (H2)
\end{verbatim}

Both scalar energy terms remain in the return. Let A\_x=H(x)-E0(x) and
{r\allowbreak{}\_\allowbreak{}p\allowbreak{}=\allowbreak{}(\allowbreak{}W\allowbreak{}\_\allowbreak{}p\allowbreak{}-\allowbreak{}\textless{}\allowbreak{}W\allowbreak{}\_\allowbreak{}p\allowbreak{}\textgreater{}\allowbreak{}\_\allowbreak{}r\allowbreak{}h\allowbreak{}o\allowbreak{})\allowbreak{}p\allowbreak{}s\allowbreak{}i\allowbreak{}\_\allowbreak{}x\allowbreak{},\allowbreak{}}
rho=psi\_x\^{}2. Define

\begin{verbatim}
C_pq(t;x)=<r_p,exp(-t A_x)r_q>,
Chat_pq(tau;x)=C_pq(tau/kappa;x),
tau=kappa t,   t=tau/kappa.                              (H3)
\end{verbatim}

This is an invertible physical-time change; no link or field coordinate
is changed. The Laplace variable z below is conjugate to this stated
tau:

\begin{verbatim}
F_pq(z;x)=int_0^infinity exp(-z tau) Chat_pq(tau;x) dtau
         =kappa <r_p,(A_x+kappa z)^(-1)r_q>.              (H4)
\end{verbatim}

All expressions are initially at real sufficiently small x and
z\textgreater0. The separate analytic proof gives their full domains and
uniform time remainders.

\subsection{H2. The complete coefficient recurrence, including the
ground
pole}\label{h2.-the-complete-coefficient-recurrence-including-the-ground-pole}

Retain the source section u=psi\_x/m\_x, m\_x=P\_H psi\_x. Its inverse
is psi\_x=m\_x u, and its original scalar obeys
m\_x\^{}2\textless u,u\textgreater\_H=1. Thus P\_H u=1 and the original
norm N(x)=\textless u,u\textgreater\_H is kept in every expectation.
Expand u(x)=sum\_nu u\_nu x\^{}nu,
{e\allowbreak{}(\allowbreak{}x\allowbreak{})\allowbreak{}=\allowbreak{}s\allowbreak{}u\allowbreak{}m\allowbreak{}\_\allowbreak{}(\allowbreak{}n\allowbreak{}u\allowbreak{}!\allowbreak{}=\allowbreak{}0\allowbreak{})\allowbreak{}e\allowbreak{}\_\allowbreak{}n\allowbreak{}u\allowbreak{}}
x\^{}nu, with u\_0=1, P\_H u\_nu=0 for nu!=0. The exact eigen-equation
gives

\begin{verbatim}
e_nu=-sum_(j:nu_j>0) P_H(W_j u_(nu-e_j)),
K u_nu=sum_j W_j u_(nu-e_j)
          +sum_(0<mu<=nu)e_mu u_(nu-mu).                 (H5)
\end{verbatim}

The last term includes e\_nu u\_0, and its Haar scalar makes the RHS
centered. Its unique centered inverse is the original K\^{}(-1)Q\_H. The
recovered coefficient engine supplies these actual finite coefficients
and all their source labels.

For a marked face q, define
{h\allowbreak{}\_\allowbreak{}q\allowbreak{}(\allowbreak{}z\allowbreak{};\allowbreak{}x\allowbreak{})\allowbreak{}=\allowbreak{}(\allowbreak{}B\allowbreak{}\_\allowbreak{}x\allowbreak{}-\allowbreak{}e\allowbreak{}(\allowbreak{}x\allowbreak{})\allowbreak{}+\allowbreak{}z\allowbreak{})\allowbreak{}\textasciicircum{}\allowbreak{}(\allowbreak{}-\allowbreak{}1\allowbreak{})\allowbreak{}W\allowbreak{}\_\allowbreak{}q\allowbreak{}}
u(x). Its coefficient equation is

\begin{verbatim}
(K+z)h_(q,nu)=W_q u_nu+sum_j W_j h_(q,nu-e_j)
                     +sum_(0<mu<=nu)e_mu h_(q,nu-mu).     (H6)
\end{verbatim}

The source index may include repeated faces; all multiplicities remain
literal multiindex coefficients. Define

\begin{verbatim}
N_nu=sum_(alpha+beta=nu)<u_alpha,u_beta>_H,
Mp_nu=sum_(alpha+beta=nu)<u_alpha,W_p u_beta>_H,
Hpq_nu=sum_(alpha+beta=nu)<u_alpha,W_p h_(q,beta)>_H.        (H7)
\end{verbatim}

Coefficient division by N, with N\_0=1, gives the original mean
mu\_p=Mp/N and the uncentered resolvent Hpq/N. The complete connected
answer is

\begin{verbatim}
F_pq=Hpq/N-mu_p mu_q/z.                                  (H8)
\end{verbatim}

The subtracted term is retained as an actual ground-state contribution.
The source shift e\_mu in (H6), both norm divisions and this term are
all required. Every resulting coefficient in the catalogue has zero
coefficient at each power z\^{}(-j), j\textgreater0. The original
positive free excitation gap also proves this cancellation analytically
near (x,z)=(0,0), after the ground projection is removed.

\subsection{H3. Exact original Casimir inversion and time
reconstruction}\label{h3.-exact-original-casimir-inversion-and-time-reconstruction}

An original invariant trace monomial has a finite tensor representation
on its original edge set. For r\_e occurrences on edge e, the allowed
doubled spins are r\_e,r\_e-2,... . At each original vertex the sum of
the doubled spins is even, and no one exceeds the sum of the others.
These are the SU(2) invariant-tensor conditions: iterated
angular-momentum addition supplies every intermediate spin in the
parity-compatible interval, so zero belongs exactly on those conditions.
The resulting finite Casimir list contains all eigenvalues of that
monomial. It retains multiplicities inside every equal-eigenvalue
eigenspace.

For its distinct values C, use the literal original operator projectors

\begin{verbatim}
P_c=product_(d in C,d!=c)(K-dI)/(c-d),
(K+z)^(-1)f=sum_(c in C) P_c f/(z+c).                     (H9)
\end{verbatim}

The full physical operator has all spins. The finite list in (H9) is the
complete representation space reached by that Taylor coefficient and its
marked word. It does not discard a path contributing to the coefficient.

Every response coefficient is stored as exact partial fractions

\begin{verbatim}
F_nu(z)=sum_(c>0,n>=1) a_(c,n)/(z+c)^n.                   (H10)
\end{verbatim}

Their time functions are exactly

\begin{verbatim}
Chat_nu(tau)=sum_(c,n) a_(c,n) tau^(n-1)exp(-c tau)/(n-1)!.(H11)
\end{verbatim}

Indeed direct integration of each finite term gives (H10). The analytic
heat construction in the companion proof shows these are the actual
Taylor coefficients, and uniqueness of the Laplace transform identifies
them. For explicit uniqueness, weighting a vanishing transform by
exp(-s0 tau) produces a finite complex measure; x=exp(-tau) carries its
integer Laplace samples to all polynomial moments on {[}0,1{]}.
Polynomial approximation forces the measure to vanish; continuity gives
equality of the time functions.

Products in the partial-fraction engine are also exact. For a!=b, the
principal part of (z+a)\^{}(-m)(z+b)\^{}(-n) at -a has coefficients

\begin{verbatim}
(-1)^k binom(n+k-1,k)/(b-a)^(n+k),  k=0,...,m-1,          (H12)
\end{verbatim}

at power (z+a)\^{}(-(m-k)); the analogous -b principal part is retained.
Their sum has the same two principal parts and vanishes at infinity, so
the rational functions are equal. Coincident poles add their orders.
This retains every repeated pole and hence every time-polynomial factor
in (H11).

The exact original source equation (H6) is audited as a complete
polynomial in all retained tree/chord quaternion coordinates, separately
for every pole. The map is Z\_e=t\_s U\_e t\_t\^{}(-1), with its
original inverse U\_e=t\_s\^{}(-1)Z\_e t\_t and the tree variables
retained. On each chord, Z=x0 I-i sum x\_a sigma\_a and sum\_(a=0)\^{}3
x\_a\^{}2=1. The remainder basis has x0 exponent zero or one. The map
and injectivity argument are the recovered fifth-source coordinate
proof, unchanged. Every residual coefficient is zero in this full
product-of-spheres algebra, not merely at sampled links.

\subsection{H4. Complete support classification and original box
assembly}\label{h4.-complete-support-classification-and-original-box-assembly}

A change U\_e -\textgreater{} -U\_e multiplies each source or marked
trace using e by -1. Thus a nonzero coefficient has an even total
multiplicity on every original edge, counting both marked traces. At
total face multiplicity 2 or 4 this requires all face multiplicities
even. At total multiplicity 6 it permits, in addition, the six distinct
faces of one elementary cube. The original closed cubical two-chain
argument is retained in the sixth-energy proof: a nonempty closed
mod-two face set has at least six faces; at six, each of the four edges
of an extreme face requires its adjacent side face, and closure of their
remaining edges forces the opposite face. The resulting set is the
boundary of that original cube.

A disconnected face source factors over disjoint original edge factors.
The original operators, positive vacua and their means factor
accordingly. A source component containing neither mark cancels between
numerator and N; marks in different components have zero connected
correlation. This remains true with shared graph vertices: the full
tensor-product Hamiltonian and the actual invariant vectors give the
same physical expectation. Connectedness therefore means sharing
original edges, exactly as in the source catalogue.

The complete list is:

{\def\LTcaptype{none} % do not increment counter
\begin{longtable}[]{@{}
  >{\raggedright\arraybackslash}p{(\linewidth - 6\tabcolsep) * \real{0.2143}}
  >{\raggedright\arraybackslash}p{(\linewidth - 6\tabcolsep) * \real{0.2143}}
  >{\raggedleft\arraybackslash}p{(\linewidth - 6\tabcolsep) * \real{0.2857}}
  >{\raggedleft\arraybackslash}p{(\linewidth - 6\tabcolsep) * \real{0.2857}}@{}}
\toprule\noalign{}
\begin{minipage}[b]{\linewidth}\raggedright
heat/source degree
\end{minipage} & \begin{minipage}[b]{\linewidth}\raggedright
total face multiplicity including marks
\end{minipage} & \begin{minipage}[b]{\linewidth}\raggedleft
coordinate cases
\end{minipage} & \begin{minipage}[b]{\linewidth}\raggedleft
marked responses
\end{minipage} \\
\midrule\noalign{}
\endhead
\bottomrule\noalign{}
\endlastfoot
0 & 2 & 1 & 1 \\
2 & 4 & 3 & 7 \\
4 & 6 & 13 & 76 \\
\end{longtable}
}

There are 84 full rational time functions. At total degree six the cases
are one single face, two geometries of an ordered 4+2 pair, seven
geometries of a three-face path, one common-edge triple, one cube corner
and one cube. All marks and all allowed diagonal marks are retained.
Faces with repeated multiplicity are never replaced by a factorial
convention.

A signed coordinate permutation and translation

\begin{verbatim}
y_i=epsilon_i x_(pi(i))-d_i,
x_(pi(i))=epsilon_i(y_i+d_i)                              (H13)
\end{verbatim}

transports every actual face multiset and both marked faces to its
record. Each original edge is mapped to the corresponding forward or
inverse edge. The product-Haar map is unitary, carries K to K on that
support, and maps the actual face words to the target face word or its
inverse. In SU(2), tr(U\^{}-1)=tr(U) by the two eigenvalues
lambda,lambda\^{}-1; hence the same trace is returned. The inverse in
(H13) returns every source to its original coordinates.

The L=2 assembly retains all 240 original faces, 1,128 adjacent pairs,
6,732 paths, 576 common-edge triples, 512 corners and 64 cubes. The
coefficient matrices have respectively 240, 2,496 and 9,660 nonzero
entries. Their entire time integrals agree entry by entry with the
preceding static matrices:

\begin{verbatim}
int_0^infinity Chat_0(tau)dtau=I/3,
int_0^infinity Chat_2(tau)dtau
  =-(5/36)I+(2/1053)D_degree+(4/1053)A_adj,                (H14)
\end{verbatim}

and the full original Chat\_4 integral is the previously calculated R4,
including every boundary row and cube. All complete matrices and all
their coordinate transports are in
{g\allowbreak{}e\allowbreak{}n\allowbreak{}e\allowbreak{}r\allowbreak{}a\allowbreak{}t\allowbreak{}e\allowbreak{}d\allowbreak{}/\allowbreak{}h\allowbreak{}e\allowbreak{}a\allowbreak{}t\allowbreak{}\_\allowbreak{}m\allowbreak{}a\allowbreak{}t\allowbreak{}r\allowbreak{}i\allowbreak{}x\allowbreak{}\_\allowbreak{}L\allowbreak{}2\allowbreak{}.\allowbreak{}j\allowbreak{}s\allowbreak{}o\allowbreak{}n\allowbreak{}}
and
{h\allowbreak{}e\allowbreak{}a\allowbreak{}t\allowbreak{}\_\allowbreak{}t\allowbreak{}r\allowbreak{}a\allowbreak{}n\allowbreak{}s\allowbreak{}p\allowbreak{}o\allowbreak{}r\allowbreak{}t\allowbreak{}s\allowbreak{}\_\allowbreak{}L\allowbreak{}2\allowbreak{}.\allowbreak{}j\allowbreak{}s\allowbreak{}o\allowbreak{}n\allowbreak{}.\allowbreak{}}
Those files preserve the original face ordering.

\subsection{H5. Explicit time-dependent
functions}\label{h5.-explicit-time-dependent-functions}

For one original plaquette with only its own source parameter, the
degree-two coefficient is

\begin{verbatim}
exp(-3tau)[-241/900-(7/15)tau]+(4/225)exp(-8tau).          (H15)
\end{verbatim}

For two different faces sharing one original edge, the coefficient of
x\_p x\_q is

\begin{verbatim}
exp(-3tau)[-409/17199+tau/21]
   +exp(-9tau/2)/36+exp(-13tau/2)/588.                    (H16)
\end{verbatim}

Its integral is 4/1053 and its value at zero is 2/351. The corresponding
diagonal, neighbour-source term replaces -409/17199 by -13/441; its
integral is 2/1053 and its time-zero value is zero. These distinct
values retain the full source and state pairings.

Now take opposite faces

\begin{verbatim}
p=(0,0,0;0,1),  q=(0,0,1;0,1).                          (H17)
\end{verbatim}

Their coefficients at degrees zero and two vanish. Four original
three-face paths and one original cube contribute at degree four. One
path gives

\begin{verbatim}
P(tau)=exp(-3tau)[75173779/67612708800
                 -(4441/3611790)tau+tau^2/882]
   -(121/108864)exp(-9tau/2)+(1/11664)exp(-6tau)
   -(197/9988160)exp(-13tau/2)+(1/163800)exp(-8tau)
   +(1/6492304)exp(-10tau).                              (H18)
\end{verbatim}

The cube gives

\begin{verbatim}
B(tau)=exp(-3tau)[-235/648+(10/27)tau]
   +exp(-9tau/2)[1123/2916+(23/81)tau+tau^2/18]
   +(1/648)exp(-6tau).                                  (H19)
\end{verbatim}

Thus the complete actual degree-four coefficient is

\begin{verbatim}
h_4(tau)=4P(tau)+B(tau),
h_4(0)=8869/365040,
int_0^infinity h_4(tau)dtau=641033/29568240.               (H20)
\end{verbatim}

All six decay rates and both repeated-pole orders remain. The
coefficient at tau=1 is enclosed by rational exponential bounds in
{g\allowbreak{}e\allowbreak{}n\allowbreak{}e\allowbreak{}r\allowbreak{}a\allowbreak{}t\allowbreak{}e\allowbreak{}d\allowbreak{}/\allowbreak{}h\allowbreak{}e\allowbreak{}a\allowbreak{}t\allowbreak{}\_\allowbreak{}b\allowbreak{}o\allowbreak{}u\allowbreak{}n\allowbreak{}d\allowbreak{}s\allowbreak{}.\allowbreak{}j\allowbreak{}s\allowbreak{}o\allowbreak{}n\allowbreak{};\allowbreak{}}
its numerical size is approximately 0.00859150755. This diagnostic
decimal is not an acceptance test. The companion proof turns (H20) into
intervals for the full actual correlation, retaining every higher Taylor
order.

\subsection{H6. Independent checks on the same original
geometry}\label{h6.-independent-checks-on-the-same-original-geometry}

The single-plaquette check uses the infinite original character basis
chi\_(j/2), where K
{c\allowbreak{}h\allowbreak{}i\allowbreak{}\_\allowbreak{}(\allowbreak{}j\allowbreak{}/\allowbreak{}2\allowbreak{})\allowbreak{}=\allowbreak{}j\allowbreak{}(\allowbreak{}j\allowbreak{}+\allowbreak{}2\allowbreak{})\allowbreak{}c\allowbreak{}h\allowbreak{}i\allowbreak{}\_\allowbreak{}(\allowbreak{}j\allowbreak{}/\allowbreak{}2\allowbreak{})\allowbreak{}}
and W
{c\allowbreak{}h\allowbreak{}i\allowbreak{}\_\allowbreak{}(\allowbreak{}j\allowbreak{}/\allowbreak{}2\allowbreak{})\allowbreak{}=\allowbreak{}c\allowbreak{}h\allowbreak{}i\allowbreak{}\_\allowbreak{}(\allowbreak{}(\allowbreak{}j\allowbreak{}-\allowbreak{}1\allowbreak{})\allowbreak{}/\allowbreak{}2\allowbreak{})\allowbreak{}+\allowbreak{}c\allowbreak{}h\allowbreak{}i\allowbreak{}\_\allowbreak{}(\allowbreak{}(\allowbreak{}j\allowbreak{}+\allowbreak{}1\allowbreak{})\allowbreak{}/\allowbreak{}2\allowbreak{})\allowbreak{},\allowbreak{}}
with the negative-index term zero. At order n every contributing
insertion path has j\textless=n+1; the independent recurrence retains
all these indices. It reproduces all three self-source heat
coefficients, with both original norm divisions and the ground pole. Its
integrated coefficients are
{1\allowbreak{}/\allowbreak{}3\allowbreak{},\allowbreak{}-\allowbreak{}5\allowbreak{}/\allowbreak{}3\allowbreak{}6\allowbreak{},\allowbreak{}2\allowbreak{}8\allowbreak{}9\allowbreak{}/\allowbreak{}5\allowbreak{}1\allowbreak{}8\allowbreak{}4\allowbreak{}.\allowbreak{}}
They also follow by twice differentiating the retained one-plaquette
energy coefficients. The original radial coordinate return
{e\allowbreak{}\_\allowbreak{}o\allowbreak{}n\allowbreak{}e\allowbreak{}(\allowbreak{}x\allowbreak{}i\allowbreak{})\allowbreak{}=\allowbreak{}b\allowbreak{}\_\allowbreak{}2\allowbreak{}(\allowbreak{}-\allowbreak{}4\allowbreak{}x\allowbreak{}i\allowbreak{})\allowbreak{}/\allowbreak{}4\allowbreak{}-\allowbreak{}1\allowbreak{}}
matches the NIST DLMF 28.6.5 expansion; the map and scalar energy term
are in the recovered audit report. This comparison covers that
one-plaquette coefficient convention only.

There is an additional computation for every marked cube pair. Orient
the six faces outward. Original Haar integration yields
2\^{}8/2\^{}12=1/16. For any proper visited face subset S, each edge
already visited twice can never occur in a later insertion. Its
surviving representation is the original spin zero; each boundary edge
carries the remaining spin one-half. Therefore its original intermediate
energy is c(S)=3\textbar boundary S\textbar/4.

For marks p!=q, arrange the other four faces in a permutation and split
the ordered list into right-vacuum, middle-heat and left-vacuum pieces.
All 24 permutations and all 15 length triples occur. Right and left
vacuum prefixes contribute 1/c(S); the middle heat state begins at \{q\}
union right and each middle insertion contributes 1/{[}z+c(S){]}. The
original 1/16 remains. Summing these 360 rational functions gives
exactly the cube coefficient in (H10). All fifteen marked cube pairs
were checked: 5,400 complete original insertion records, including every
intermediate boundary energy. This calculation imports none of the
trace-differentiation or Casimir-projection engine.

The time-zero covariance, its first derivative
{-\allowbreak{}\textless{}\allowbreak{}G\allowbreak{}a\allowbreak{}m\allowbreak{}m\allowbreak{}a\allowbreak{}(\allowbreak{}W\allowbreak{}\_\allowbreak{}p\allowbreak{},\allowbreak{}W\allowbreak{}\_\allowbreak{}q\allowbreak{})\allowbreak{}\textgreater{}\allowbreak{}\_\allowbreak{}r\allowbreak{}h\allowbreak{}o\allowbreak{},\allowbreak{}}
and its second derivative
{\textless{}\allowbreak{}[\allowbreak{}K\allowbreak{},\allowbreak{}W\allowbreak{}\_\allowbreak{}p\allowbreak{}]\allowbreak{}u\allowbreak{},\allowbreak{}[\allowbreak{}K\allowbreak{},\allowbreak{}W\allowbreak{}\_\allowbreak{}q\allowbreak{}]\allowbreak{}u\allowbreak{}\textgreater{}\allowbreak{}\_\allowbreak{}H\allowbreak{}/\allowbreak{}N\allowbreak{}}
are also calculated separately from the entire original density quotient
at each coefficient. All 84 rational functions reproduce these three
values. The full-polynomial source audit and these moment identities
have distinct execution records.

\subsection{H7. Sources and present mathematical
scope}\label{h7.-sources-and-present-mathematical-scope}

The original operator/domain and maximal-tree maps are the pinned
finite-box workbench source. The exponential and character/Casimir
antecedent is D. Schütte, Zheng Weihong and C. J. Hamer,
{a\allowbreak{}r\allowbreak{}X\allowbreak{}i\allowbreak{}v\allowbreak{}:\allowbreak{}h\allowbreak{}e\allowbreak{}p\allowbreak{}-\allowbreak{}l\allowbreak{}a\allowbreak{}t\allowbreak{}/\allowbreak{}9\allowbreak{}6\allowbreak{}0\allowbreak{}3\allowbreak{}0\allowbreak{}2\allowbreak{}6\allowbreak{}.\allowbreak{}}
The residue and minimum-section comparison imported from the current
Split-Zero heat work is specified completely in
{R\allowbreak{}H\allowbreak{}\_\allowbreak{}H\allowbreak{}E\allowbreak{}A\allowbreak{}T\allowbreak{}\_\allowbreak{}T\allowbreak{}R\allowbreak{}A\allowbreak{}N\allowbreak{}S\allowbreak{}F\allowbreak{}E\allowbreak{}R\allowbreak{}.\allowbreak{}m\allowbreak{}d\allowbreak{},\allowbreak{}}
with its exact source revision and read equations. No arithmetic period
or Gamma constant is assigned to a physical Yang--Mills matrix.

These are actual finite-regulator heat coefficients and actual
finite-volume analytic bounds. The Cauchy radius in the companion is
1/M. It therefore retains dependence on the original volume, and does
not supply an ultraviolet continuum construction. The earlier
uniform-gap manuscripts retain their recorded independent-audit
qualification.

\subsection{H8. The complete fourth-order first-band return in its
original
metric}\label{h8.-the-complete-fourth-order-first-band-return-in-its-original-metric}

The original free energy-three space is exactly span\{Wp\}. To see the
next separation, an active invariant graph with four edges is a
four-cycle; its vertex intertwiners force equal edge spins, giving
energy 4j(j+1)=3,8,... . A simple bipartite active graph with five edges
and minimum degree at least two would have one cyclic component with
five vertices and five edges, hence an odd five-cycle, or at most four
vertices, where bipartiteness permits at most four edges. Both cases
contradict its assumptions. Therefore any other active assignment has at
least six edges and energy \textgreater=9/2. Thus

\begin{verbatim}
spec(K_phys) subset {0,3} union [9/2,infinity),
P3H_phys=span{Wp},   dim P3=M.                           (H21)
\end{verbatim}

For sufficiently small real xi let Pi\_xi be the complete spectral
projection near energy 3 of B\_xi-e(xi), and put

\begin{verbatim}
Y_xi c=Pi_xi sum_p c_p r_p,
Gband=Y_xi*Y_xi,
(B_xi-e)Y_xi=Y_xi Aband.                                (H22)
\end{verbatim}

The exact quantitative construction below proves invertibility of this
frame onto the actual band. At xi=0, c-\textgreater sum c\_p Wp is the
original Haar isometry R0. Projection of the heat resolvent onto its
original energy-three poles gives

\begin{verbatim}
Cband(tau)=Gband exp(-tau Aband),
Gband=I+xi^2 G2+xi^4 G4+...,
Aband=3I+xi^2 T2+xi^4 T4+....                            (H23)
\end{verbatim}

The original center symmetry makes both matrices even. For each of the
full heat matrices let A\_(n,j) be the coefficient of (z+3)\^{}(-j).
Then literal multiplication of the two series in (H23) proves

\begin{verbatim}
G2=A_(2,1),  G4=A_(4,1),  T2=-A_(2,2),
A_(4,3)=T2^2,
T4=-A_(4,2)-G2 T2.                                    (H24)
\end{verbatim}

All matrix factors retain their order. The complete L=2 calculation
gives

\begin{verbatim}
T2=(7/15)I-(D_degree+A_adj)/21.                          (H25)
\end{verbatim}

It gives all 9,660 nonzero entries of T4 and all 9,660 entries of G4 in
{f\allowbreak{}i\allowbreak{}r\allowbreak{}s\allowbreak{}t\allowbreak{}\_\allowbreak{}b\allowbreak{}a\allowbreak{}n\allowbreak{}d\allowbreak{}\_\allowbreak{}L\allowbreak{}2\allowbreak{}.\allowbreak{}j\allowbreak{}s\allowbreak{}o\allowbreak{}n\allowbreak{}.\allowbreak{}}
Every entry of A\_(4,3)=T2\^{}2 is checked. Original self-adjointness
returns through the full Gram as

\begin{verbatim}
Gband Aband=Aband* Gband,
T4-T4* = T2 G2-G2 T2.                                  (H26)
\end{verbatim}

The right side has 1,440 nonzero entries for this open box. For the
original faces (-2,-2,-2;0,1) and (-2,-1,-2;0,1), the displayed
difference is exactly 2/7371. This is the explicit correction between
the coordinate adjoint and the original metric adjoint. Its value is
retained, rather than forcing a symmetric coefficient matrix by dropping
the mixed product.

Here is a complete analytic remainder for (H23). Work on the complex
homogeneous circle \textbar zeta\textbar=1/(32M). The original
ground/complement calculation in A1-\/-A2 applies with
{|\allowbreak{}|\allowbreak{}W\allowbreak{}|\allowbreak{}|\allowbreak{}\textasciicircum{}\allowbreak{}2\allowbreak{}=\allowbreak{}e\allowbreak{}t\allowbreak{}a\allowbreak{}=\allowbreak{}1\allowbreak{}/\allowbreak{}(\allowbreak{}1\allowbreak{}0\allowbreak{}2\allowbreak{}4\allowbreak{}M\allowbreak{})\allowbreak{}}
and Re D\textgreater=47/16. On \textbar e\textbar=eta,
{|\allowbreak{}|\allowbreak{}Z\allowbreak{}(\allowbreak{}D\allowbreak{}-\allowbreak{}e\allowbreak{})\allowbreak{}\textasciicircum{}\allowbreak{}(\allowbreak{}-\allowbreak{}1\allowbreak{})\allowbreak{}W\allowbreak{}|\allowbreak{}|\allowbreak{}\textless{}\allowbreak{}e\allowbreak{}t\allowbreak{}a\allowbreak{},\allowbreak{}}
so the same scalar argument principle gives
\textbar e\textbar\textless=eta. Consequently

\begin{verbatim}
||B_zeta-e-K|| <=1/16+1/(1024M)<1/15.                    (H27)
\end{verbatim}

On the spectral circle \textbar lambda-3\textbar=1/2 the original free
resolvent has norm at most 2. Its norm-convergent resolvent series gives

\begin{verbatim}
||Pi-P3|| <=2/13,
||(B_zeta-e-3)Pi||<=1/13.                               (H28)
\end{verbatim}

For the first inequality, integrate the difference of resolvents: the
bound is 2b/(1-2b) with b\textless=1/15. For the second, integrate that
difference multiplied by lambda-3, obtaining b/(1-2b). The corresponding
free integral is zero because (K-3)P3=0. The circle encloses exactly the
original M-dimensional band, by the same resolvent homotopy and finite
rank at zero.

Use the exact Haar section u=1+w, where
{|\allowbreak{}|\allowbreak{}w\allowbreak{}|\allowbreak{}|\allowbreak{}\textless{}\allowbreak{}=\allowbreak{}1\allowbreak{}/\allowbreak{}(\allowbreak{}9\allowbreak{}2\allowbreak{}s\allowbreak{}q\allowbreak{}r\allowbreak{}t\allowbreak{}(\allowbreak{}M\allowbreak{})\allowbreak{})\allowbreak{};\allowbreak{}}
the denominator in its block inverse exceeds 23/8. The column map
Fu:c-\textgreater sum c\_p Wp u satisfies

\begin{verbatim}
||Fu-R0||<=2sqrt(M)||w||<=1/46.
\end{verbatim}

Define Yhat=Pi Fu. Its original coefficient map obeys

\begin{verbatim}
||R0*Yhat-I||<=107/598,
||(R0*Yhat)^(-1)||<=598/491.                             (H29)
\end{verbatim}

Indeed the two contributions are Pi-P3 and Pi(Fu-R0), bounded by 2/13
and (15/13)/46. Its inverse on the band is (R0*Yhat)\^{}(-1)R0*. The
actual unit-vacuum frame is Y=m Yhat, with psi=m u and
m\^{}2\textless u,u\textgreater\_H=1 retained. Its scalar cancels from
the operator intertwining, not from Gband. Thus

\begin{verbatim}
Aband=(R0*Yhat)^(-1)R0*(B_zeta-e)Yhat,
||Aband-3I|| <=(598/491)(1/13)(47/46)=47/491.             (H30)
\end{verbatim}

All these maps are analytic on a neighborhood of the closed source
circle. The original center involution sends Yhat\_-zeta=-C Yhat\_zeta
and R0*C=-R0*, so Aband is even. Cauchy\textquotesingle s formula and
the full even tail prove

\begin{verbatim}
||Aband-3I-xi^2T2-xi^4T4||
  <=(47/491)(32M|xi|)^6/[1-(32M|xi|)^2],
                     |xi|<1/(32M).                    (H31)
\end{verbatim}

Multiplication by kappa returns this to physical energy. This is an
actual finite-volume matrix remainder, with all original band
coordinates, the nonidentity Gram (H23), and the dependence on M
retained. No diagonalization of T4 or uniform-in-volume fourth-order
spectral claim is inferred here.

\section{Full-time heat bounds and the original plaquette response
metrics}\label{full-time-heat-bounds-and-the-original-plaquette-response-metrics}

21 September 2026. This proof uses the unchanged finite-box SU(2)
Hamiltonian and the complete coefficients in
{P\allowbreak{}H\allowbreak{}Y\allowbreak{}S\allowbreak{}I\allowbreak{}C\allowbreak{}A\allowbreak{}L\allowbreak{}\_\allowbreak{}H\allowbreak{}E\allowbreak{}A\allowbreak{}T\allowbreak{}\_\allowbreak{}C\allowbreak{}O\allowbreak{}E\allowbreak{}F\allowbreak{}F\allowbreak{}I\allowbreak{}C\allowbreak{}I\allowbreak{}E\allowbreak{}N\allowbreak{}T\allowbreak{}S\allowbreak{}.\allowbreak{}m\allowbreak{}d\allowbreak{}.\allowbreak{}}
Its analytic argument is finite-volume and independent of the inherited
volume-uniform source-norm estimates. The original number M of
plaquettes appears in every Cauchy radius. Written arguments and
executed finite checks have separate records. All inner products are
conjugate-linear in the first entry.

\subsection{A1. Original free spectrum and Haar
blocks}\label{a1.-original-free-spectrum-and-haar-blocks}

Use H1-\/-H3 of the coefficient proof. The original physical free
operator K has its constant eigenvector 1, and every other physical
Fourier block has energy at least 3. Here is the graph argument. In a
nonzero invariant representation assignment, a vertex incident to only
one active edge has no invariant tensor. Thus the finite active graph
has minimum degree at least two and contains a cycle. The original
simple cubic graph has no triangles and every cycle has at least four
edges. Each active edge has spin j\textgreater=1/2 and Casimir
j(j+1)\textgreater=3/4. Hence the complete physical K on the orthogonal
complement of 1 is \textgreater=3. The elementary plaquette trace has
four spin-one-half edges and energy 3.

Write P0=\textbar1\textgreater\textless1\textbar{} and Q0=I-P0 in the
original product Haar space. Integrating one original link gives P0
Wp=0. Two distinct elementary plaquettes have an edge occurring in just
one of their two words; integration over that link vanishes. For p=q,
the product holonomy has the original Haar law and the fundamental
character has squared norm one. Therefore

\begin{verbatim}
<Wp,Wq>_H=delta_pq.                                      (A1)
\end{verbatim}

For the complex homogeneous coefficient zeta, x\_p=zeta, put
B\_zeta=K-zeta S. This complex variable retains the physical return H=
kappa B\_xi+2kappa M xi I at real xi. On C1 plus Q0 H\_phys the full
original block operator is

\begin{verbatim}
B_zeta = [[0,Z],[W,D]],
D=Q0(K-zeta S)Q0 on Dom(K) intersect Q0 H_phys,
W=-zeta sum_p Wp,   Z=-zeta P0 S Q0.                     (A2)
\end{verbatim}

Z is the original complex-linear row, with no conjugation of its
parameter in the analytic extension. Equation (A1) proves exactly

\begin{verbatim}
||W||=||Z||=sqrt(M)|zeta|.
\end{verbatim}

On \textbar zeta\textbar\textless=1/M, Re\textless Dh,h\textgreater{}
\textgreater= \textbar\textbar h\textbar\textbar\^{}2 since
\textbar\textbar zeta S\textbar\textbar\textless=2. The closed operator
D is a bounded perturbation of Q0KQ0 on the same domain. For Re
e\textless1 its inverse D-e exists with norm \textless=(1-Re e)\^{}(-1):
injectivity and closed range follow from the real part inequality; the
same inequality for its adjoint gives dense range, hence surjectivity.
All finite boxes here have
{M\allowbreak{}=\allowbreak{}3\allowbreak{}(\allowbreak{}2\allowbreak{}L\allowbreak{})\allowbreak{}\textasciicircum{}\allowbreak{}2\allowbreak{}(\allowbreak{}2\allowbreak{}L\allowbreak{}+\allowbreak{}1\allowbreak{})\allowbreak{}\textgreater{}\allowbreak{}=\allowbreak{}2\allowbreak{}4\allowbreak{}0\allowbreak{}.\allowbreak{}}

\subsection{A2. A complete ground/complement decomposition on the
complex
disk}\label{a2.-a-complete-groundcomplement-decomposition-on-the-complex-disk}

The scalar equation for the small eigenvalue is

\begin{verbatim}
f(e)=e+Z(D-e)^(-1)W=0.                                  (A3)
\end{verbatim}

On \textbar e\textbar=2/M the second term has modulus at most
1/(M-2)\textless2/M. The finite scalar argument principle, or
Rouche\textquotesingle s theorem applied to e, gives one simple zero
inside that circle, with multiplicity counted. The contour construction
makes e(zeta) analytic on a neighborhood of the closed
\textbar zeta\textbar\textless=1/M disk and
\textbar e\textbar\textless=2/M. Define the actual right and left graph
coordinates

\begin{verbatim}
w=-(D-e)^(-1)W,
ell=-Z(D-e)^(-1),
P=(1,w)(1,ell)/(1+ell w).                               (A4)
\end{verbatim}

They satisfy B(1,w)=e(1,w), (1,ell)B=e(1,ell), ell W=e, and Z+ell D=e
ell. Their norms obey

\begin{verbatim}
||w||,||ell|| <= sqrt(M)/(M-2),
M/(M-2)^2 <= 240/238^2 < 1/225.
\end{verbatim}

The last bound decreases with M\textgreater=240 by differentiating
M/(M-2)\^{}2. Consequently 1+ell w is nonzero, P\^{}2=P, and

\begin{verbatim}
||P||_trace <= (1+1/225)/(1-1/225)=113/112.               (A5)
\end{verbatim}

The original complement is ker(1,ell). Its exact coordinate map and
inverse are

\begin{verbatim}
J:Q0H -> ker(1,ell),   Jh=(-ell h,h),
J^(-1)=Q0 restricted to ker(1,ell).                     (A6)
\end{verbatim}

On the unchanged operator domain direct multiplication gives

\begin{verbatim}
(B-e)J=J Bc,   Bc=D-e-Well.                             (A7)
\end{verbatim}

Every factor in the feedback Well is retained. For its real part,

\begin{verbatim}
Re<Bc h,h> >= [1-2/M-1/(M-2)]||h||^2
             >=[1-1/120-1/238]||h||^2
             > (49/50)||h||^2.                         (A8)
\end{verbatim}

Bc is again the same self-adjoint free operator plus a bounded
perturbation. The norm-convergent Duhamel series constructs its strongly
continuous semigroup; its derivative on Dom(K), followed by the
real-part energy inequality (A8), proves \textbar\textbar exp(-tau
{B\allowbreak{}c\allowbreak{})\allowbreak{}|\allowbreak{}|\allowbreak{}\textless{}\allowbreak{}=\allowbreak{}e\allowbreak{}x\allowbreak{}p\allowbreak{}(\allowbreak{}-\allowbreak{}4\allowbreak{}9\allowbreak{}t\allowbreak{}a\allowbreak{}u\allowbreak{}/\allowbreak{}5\allowbreak{}0\allowbreak{})\allowbreak{}.\allowbreak{}}
Density extends this estimate to all vectors. Holomorphy in zeta follows
locally from the uniformly convergent bounded perturbation series and
the analytic coefficients in (A4).

The exact excited semigroup is therefore

\begin{verbatim}
exp[-tau(B-e)](I-P)=J exp(-tau Bc)Q0(I-P).               (A9)
\end{verbatim}

The bounds on the original graph maps are

\begin{verbatim}
||J||<=226/225,
||Q0P||<=113/1680,
||Q0(I-P)||<=1793/1680.                                 (A10)
\end{verbatim}

For example
{Q\allowbreak{}0\allowbreak{}P\allowbreak{}=\allowbreak{}w\allowbreak{}(\allowbreak{}1\allowbreak{},\allowbreak{}e\allowbreak{}l\allowbreak{}l\allowbreak{})\allowbreak{}/\allowbreak{}(\allowbreak{}1\allowbreak{}+\allowbreak{}e\allowbreak{}l\allowbreak{}l\allowbreak{}}
w), and using \textbar\textbar w\textbar\textbar\textless=1/15 and
{s\allowbreak{}q\allowbreak{}r\allowbreak{}t\allowbreak{}(\allowbreak{}1\allowbreak{}+\allowbreak{}1\allowbreak{}/\allowbreak{}2\allowbreak{}2\allowbreak{}5\allowbreak{})\allowbreak{}\textless{}\allowbreak{}=\allowbreak{}2\allowbreak{}2\allowbreak{}6\allowbreak{}/\allowbreak{}2\allowbreak{}2\allowbreak{}5\allowbreak{}}
gives 113/1680. No product-space Gram is replaced.

For real xi, B\_xi is the full self-adjoint physical operator less its
displayed scalar. Its lowest eigenvalue is \textless=0 by the original
constant trial state. Equation (A8) leaves the entire remaining spectrum
positive after subtracting e. Thus e is its actual ground eigenvalue. P
is its actual orthogonal ground projection, obtained above on the same
domain. This also identifies the analytic branch without choosing a
surrogate vacuum.

\subsection{A3. The complete all-time analytic
bound}\label{a3.-the-complete-all-time-analytic-bound}

Define the analytic connected heat observable by the literal trace

\begin{verbatim}
Chat_pq(tau;zeta)
  =Tr[P Wp exp[-tau(B-e)](I-P)Wq].                       (A11)
\end{verbatim}

At real xi this equals the original
{C\allowbreak{}\_\allowbreak{}p\allowbreak{}q\allowbreak{}(\allowbreak{}t\allowbreak{}a\allowbreak{}u\allowbreak{}/\allowbreak{}k\allowbreak{}a\allowbreak{}p\allowbreak{}p\allowbreak{}a\allowbreak{};\allowbreak{}x\allowbreak{}i\allowbreak{})\allowbreak{},\allowbreak{}}
with its actual vacuum-mean subtraction. The complex trace is its
holomorphic continuation; no conjugation of zeta is inserted. Each
original multiplication Wp has norm at most 2. Equations (A5), (A9),
(A10) give

\begin{verbatim}
|Chat_pq(tau;zeta)|
  <=4(113/112)(226/225)(1793/1680)exp(-49tau/50)
  <5 exp(-49tau/50),    |zeta|<=1/M, tau>=0.             (A12)
\end{verbatim}

The original link-center substitution U\_(n,i) -\textgreater{}
{(\allowbreak{}-\allowbreak{}1\allowbreak{})\allowbreak{}\textasciicircum{}\allowbreak{}(\allowbreak{}s\allowbreak{}u\allowbreak{}m\allowbreak{}\_\allowbreak{}(\allowbreak{}j\allowbreak{}\textless{}\allowbreak{}i\allowbreak{})\allowbreak{}n\allowbreak{}\_\allowbreak{}j\allowbreak{})\allowbreak{}}
U\_(n,i) sends every plaquette trace to -Wp, preserves Haar, K and all
gauge spaces, and is its own inverse. It carries B\_zeta to B\_-zeta.
Both marked traces change sign, so (A11) is even in zeta.
Cauchy\textquotesingle s coefficient formula on the original circle
\textbar zeta\textbar=1/M and summation of the full even tail now prove

\begin{verbatim}
|Chat_pq(tau;xi)-C0_pq(tau)-xi^2 C2_pq(tau)-xi^4 C4_pq(tau)|
  <=5 exp(-49tau/50)(M|xi|)^6/[1-(M|xi|)^2],
                |xi|<1/M, tau>=0.                     (A13)
\end{verbatim}

All constants and all higher orders remain. The bound is uniform in
positive time, with its original volume dependence. It also justifies
differentiation in the source before Laplace integration. The finite
rational functions in the coefficient proof are therefore precisely the
Taylor coefficients of (A11). The Laplace uniqueness argument H3 returns
their actual time functions.

\subsection{A4. A positive original opposite-face correlation on an
entire time
interval}\label{a4.-a-positive-original-opposite-face-correlation-on-an-entire-time-interval}

For the two faces H17, C0=C2=0 and C4=h4=4P+B with H18-\/-H20. Exact
rational Taylor remainder estimates enclose exp(-x) for each x in
{[}0,50{]}: the even Taylor sum of degree 240 is an upper bound, and
subtracting x\^{}241/241! is a lower bound. Taylor\textquotesingle s
Lagrange remainder has the required negative sign and magnitude
\textless=x\^{}241/241!. The code checks positive lower endpoints. This
is an interval calculation with rational endpoints, including at all
cited times.

At L=2, M=240, xi=10\^{}(-10), tau=1, evaluation of H18-\/-H20 and (A13)
gives

\begin{verbatim}
85879/10^7 < Chat_pq(1;xi)/xi^4 < 85951/10^7.             (A14)
\end{verbatim}

The complete rational inner endpoints are in
{g\allowbreak{}e\allowbreak{}n\allowbreak{}e\allowbreak{}r\allowbreak{}a\allowbreak{}t\allowbreak{}e\allowbreak{}d\allowbreak{}/\allowbreak{}h\allowbreak{}e\allowbreak{}a\allowbreak{}t\allowbreak{}\_\allowbreak{}b\allowbreak{}o\allowbreak{}u\allowbreak{}n\allowbreak{}d\allowbreak{}s\allowbreak{}.\allowbreak{}j\allowbreak{}s\allowbreak{}o\allowbreak{}n\allowbreak{}.\allowbreak{}}
The physical parameters are g\^{}2=50000, kappa=100000/a, and t=1/kappa.

The positivity on 1\textless=tau\textless=3 has a continuous proof. In
the first bracket of H18, let a=1/882, b=-4441/3611790,
{c\allowbreak{}=\allowbreak{}7\allowbreak{}5\allowbreak{}1\allowbreak{}7\allowbreak{}3\allowbreak{}7\allowbreak{}7\allowbreak{}9\allowbreak{}/\allowbreak{}6\allowbreak{}7\allowbreak{}6\allowbreak{}1\allowbreak{}2\allowbreak{}7\allowbreak{}0\allowbreak{}8\allowbreak{}8\allowbreak{}0\allowbreak{}0\allowbreak{}.\allowbreak{}}
Its quadratic minimum is c-b\^{}2/(4a). Exact rational comparison gives

\begin{verbatim}
c-b^2/(4a) > (121/108864)/4+(197/9988160)/16.              (A15)
\end{verbatim}

Since exp(-3/2)\textless1/4 and exp(-7/2)\textless1/16, the two negative
exponentials in H18 are bounded by this quantity times exp(-3tau) for
tau\textgreater=1. Its other terms are positive. Thus
P(tau)\textgreater0 throughout tau\textgreater=1. In H19 the first
bracket is at least 5/648 at tau=1 and increases, and both other terms
are positive. Hence

\begin{verbatim}
h4(tau) >= (5/648)exp(-3tau),    tau>=1.                 (A16)
\end{verbatim}

For tau in {[}1,3{]}, the ratio of the complete error (A13) to xi\^{}4
times (A16) is at most

\begin{verbatim}
648 M^6 xi^2 exp(303/50)/[1-(Mxi)^2] < 0.531 < 1         (A17)
\end{verbatim}

at the actual M and xi above. Consequently the full interacting-vacuum
correlation, not just its leading coefficient, is strictly positive on
{1\allowbreak{}/\allowbreak{}k\allowbreak{}a\allowbreak{}p\allowbreak{}p\allowbreak{}a\allowbreak{}\textless{}\allowbreak{}=\allowbreak{}t\allowbreak{}\textless{}\allowbreak{}=\allowbreak{}3\allowbreak{}/\allowbreak{}k\allowbreak{}a\allowbreak{}p\allowbreak{}p\allowbreak{}a\allowbreak{}.\allowbreak{}}
Its original integrated value has all physical factors retained by H4.

\subsection{A5. Full inverse-power metrics, including the original state
norm}\label{a5.-full-inverse-power-metrics-including-the-original-state-norm}

Let R:C\^{}M-\textgreater H\_phys,0 have the original centered columns
r\_p. For integer k\textgreater=1,

\begin{verbatim}
G^(k)=kappa^k R* A^(-k) R
     =int_0^infinity tau^(k-1) Chat(tau)d tau/(k-1)!.     (A18)
\end{verbatim}

The actual finite-regulator spectral theorem and the positive centered
gap established in A2 justify the integral. Equation (A12) gives a
uniform integrable majorant for the analytic Taylor coefficients and
their remainder. For a partial fraction a/(z+c)\^{}n, its exact
contribution is

\begin{verbatim}
a (k+n-2)!/[(k-1)!(n-1)! c^(k+n-1)].                    (A19)
\end{verbatim}

Every original pole and multiplicity is used. Integrating (A13) proves
the entrywise bound

\begin{verbatim}
|G^(k)_pq - sum_(j=0)^2 xi^(2j) G^(k)_(2j),pq|
  <=5(50/49)^k (M|xi|)^6/[1-(M|xi|)^2].                 (A20)
\end{verbatim}

The physical energy response is G\^{}(1). For the actual derivative
primitive

\begin{verbatim}
Phi=kappa A^(-1)R,
A Phi=kappa R,    Phi*Phi=G^(2),
Phi*A Phi=kappa G^(1),                                 (A21)
\end{verbatim}

the original state and energy metrics remain separately specified.

All coefficients for k=1,2 on the original 240-face L=2 box are
calculated entry by entry in
{g\allowbreak{}e\allowbreak{}n\allowbreak{}e\allowbreak{}r\allowbreak{}a\allowbreak{}t\allowbreak{}e\allowbreak{}d\allowbreak{}/\allowbreak{}n\allowbreak{}a\allowbreak{}t\allowbreak{}i\allowbreak{}v\allowbreak{}e\allowbreak{}\_\allowbreak{}i\allowbreak{}n\allowbreak{}v\allowbreak{}e\allowbreak{}r\allowbreak{}s\allowbreak{}e\allowbreak{}\_\allowbreak{}m\allowbreak{}e\allowbreak{}t\allowbreak{}r\allowbreak{}i\allowbreak{}c\allowbreak{}s\allowbreak{}\_\allowbreak{}L\allowbreak{}2\allowbreak{}.\allowbreak{}j\allowbreak{}s\allowbreak{}o\allowbreak{}n\allowbreak{}.\allowbreak{}}
Let L\_(k,j) be the maximum absolute row sum of the exact degree-2j
coefficient matrix. A real symmetric M by M error whose entries have
magnitude at most e has operator norm at most M e: apply the row and
column bounds and Cauchy-\/-Schwarz. Therefore (A20), with the complete
off-diagonal rows, gives

\begin{verbatim}
[3^(-k)-w_k]I <= G^(k) <= [3^(-k)+w_k]I,
w_k=xi^2 L_(k,1)+xi^4 L_(k,2)
      +5M(50/49)^k(Mxi)^6/[1-(Mxi)^2].                  (A22)
\end{verbatim}

For 0\textless xi\textless=1/1600, all the scalar terms on the right
increase with xi, so its endpoint evaluation bounds the whole interval.
Exact rational evaluation gives

\begin{verbatim}
(3/10)I_240 < G^(1) < (11/30)I_240,
(9/100)I_240 < G^(2) < (13/100)I_240.                    (A23)
\end{verbatim}

This is the original L=2 box, every a\textgreater0, and every
g\^{}2\textgreater=20. The identities use the original coefficient
pairing on C\^{}240; the actual state Gram is the calculated G\^{}(2),
not the identity. The complete original energy form on these primitive
columns is kappa G\^{}(1). The tighter rational endpoints and each
original row-sum certificate are retained in
{g\allowbreak{}e\allowbreak{}n\allowbreak{}e\allowbreak{}r\allowbreak{}a\allowbreak{}t\allowbreak{}e\allowbreak{}d\allowbreak{}/\allowbreak{}h\allowbreak{}e\allowbreak{}a\allowbreak{}t\allowbreak{}\_\allowbreak{}b\allowbreak{}o\allowbreak{}u\allowbreak{}n\allowbreak{}d\allowbreak{}s\allowbreak{}.\allowbreak{}j\allowbreak{}s\allowbreak{}o\allowbreak{}n\allowbreak{}.\allowbreak{}}

\subsection{A6. Scope and comparison with the preceding
analysis}\label{a6.-scope-and-comparison-with-the-preceding-analysis}

The finite-volume Cauchy circle in (A13) has radius 1/M. Equation (A23)
has been evaluated on the original L=2 family, with all 240 columns.
Every dependence on M, a, g and time is displayed. These statements
establish neither a nontrivial four-dimensional continuum field nor a
uniform-in-volume positive mass bound. They do not independently
recertify the earlier
{u\allowbreak{}n\allowbreak{}i\allowbreak{}f\allowbreak{}o\allowbreak{}r\allowbreak{}m\allowbreak{}-\allowbreak{}s\allowbreak{}o\allowbreak{}u\allowbreak{}r\allowbreak{}c\allowbreak{}e\allowbreak{}/\allowbreak{}g\allowbreak{}a\allowbreak{}p\allowbreak{}}
argument.

The method uses classical bounded perturbation, block resolvents, the
scalar argument principle and the original spectral theorem. A relevant
primary Feshbach-\/-Schur antecedent is Dusson-\/-Sigal-\/-Stamm,
arXiv:2105.02058. The present proof gives all blocks, graph inverses and
finite constants on the actual Hamiltonian instead of supplying a
desired spectral gap as an input. The new Split-Zero heat transfer to
the resulting actual Grams is in
{R\allowbreak{}H\allowbreak{}\_\allowbreak{}H\allowbreak{}E\allowbreak{}A\allowbreak{}T\allowbreak{}\_\allowbreak{}T\allowbreak{}R\allowbreak{}A\allowbreak{}N\allowbreak{}S\allowbreak{}F\allowbreak{}E\allowbreak{}R\allowbreak{}.\allowbreak{}m\allowbreak{}d\allowbreak{}.\allowbreak{}}

\subsection{A7. The original forcing Gram for the same observation
family}\label{a7.-the-original-forcing-gram-for-the-same-observation-family}

The time-zero matrix is the original forcing Gram G\^{}(0)=R*R=Chat(0).
It is computed from exactly the same complete pole coefficients by

\begin{verbatim}
G^(0)_(2j),pq=sum_(c>0) a_(c,1),                        (A24)
\end{verbatim}

because higher pole orders vanish at tau=0. Equation (A13) at zero
supplies its complete entry error
{5\allowbreak{}(\allowbreak{}M\allowbreak{}x\allowbreak{}i\allowbreak{})\allowbreak{}\textasciicircum{}\allowbreak{}6\allowbreak{}/\allowbreak{}[\allowbreak{}1\allowbreak{}-\allowbreak{}(\allowbreak{}M\allowbreak{}x\allowbreak{}i\allowbreak{})\allowbreak{}\textasciicircum{}\allowbreak{}2\allowbreak{}]\allowbreak{}.\allowbreak{}}
The identical full row-sum argument on the actual L=2 box and
0\textless xi\textless=1/1600 proves

\begin{verbatim}
(49/50)I_240 < G^(0) < (51/50)I_240.                    (A25)
\end{verbatim}

No division by a chosen source norm is involved. The tighter rational
interval is retained with G\^{}(1),G\^{}(2) in
{g\allowbreak{}e\allowbreak{}n\allowbreak{}e\allowbreak{}r\allowbreak{}a\allowbreak{}t\allowbreak{}e\allowbreak{}d\allowbreak{}/\allowbreak{}h\allowbreak{}e\allowbreak{}a\allowbreak{}t\allowbreak{}\_\allowbreak{}b\allowbreak{}o\allowbreak{}u\allowbreak{}n\allowbreak{}d\allowbreak{}s\allowbreak{}.\allowbreak{}j\allowbreak{}s\allowbreak{}o\allowbreak{}n\allowbreak{}.\allowbreak{}}
Thus all three original Grams R*R, Phi*A Phi/kappa and Phi*Phi are
quantitatively controlled together. Their use in the actual
inverse-power observation is proved in T7 of the source-transfer note.

\subsection{A8. The full original kinetic Gram, retaining its local zero
entries}\label{a8.-the-full-original-kinetic-gram-retaining-its-local-zero-entries}

The forcing columns have original energy Gram R*A R=kappa Kobs, where

\begin{verbatim}
Kobs_pq=<Gamma(Wp,Wq)>_rho=-Chat'_pq(0).                 (A26)
\end{verbatim}

The ground-state form identity proves the first equality;
differentiation of the original semigroup proves the second. In the
coefficient file each simple pole contributes its energy times its
coefficient and each double pole contributes minus its coefficient.
Higher poles contribute zero. This gives all entries through degree four
in
{n\allowbreak{}a\allowbreak{}t\allowbreak{}i\allowbreak{}v\allowbreak{}e\allowbreak{}\_\allowbreak{}i\allowbreak{}n\allowbreak{}v\allowbreak{}e\allowbreak{}r\allowbreak{}s\allowbreak{}e\allowbreak{}\_\allowbreak{}m\allowbreak{}e\allowbreak{}t\allowbreak{}r\allowbreak{}i\allowbreak{}c\allowbreak{}s\allowbreak{}\_\allowbreak{}L\allowbreak{}2\allowbreak{}.\allowbreak{}j\allowbreak{}s\allowbreak{}o\allowbreak{}n\allowbreak{}.\allowbreak{}}

Every off-diagonal entry outside the original face adjacency is zero
exactly: Gamma differentiates a common original edge. On a diagonal,
Gamma(Wp,Wp)=4-Wp\^{}2 has absolute value at most4. On adjacent distinct
faces, their single shared-edge derivative vectors each have Euclidean
length at most one; writing the original SU(2) product as q0I-iq.sigma
gives the length sqrt(1-q0\^{}2). Their scalar product has absolute
value at most1. Each face has at most12 original neighbours. All
assertions concern the unchanged original link derivatives, with their
generator factor1/2 retained.

The complex ground projection from (A4)-\/-(A5) therefore bounds the
diagonal expectation by4*113/112 and each adjacent expectation
by113/112. Original center symmetry makes these analytic expectations
even. Cauchy\textquotesingle s formula on \textbar zeta\textbar=1/M
bounds the entire row of the omitted tail by

\begin{verbatim}
(113/112)*16*(M|xi|)^6/[1-(M|xi|)^2].                   (A27)
\end{verbatim}

This uses the exact local support, rather than a dense-M multiplier on
this particular Gram. The exact degree-two and degree-four row sums and
(A27) give, on the original L=2 box with g\^{}2\textgreater=20,

\begin{verbatim}
(29/10)I_240 < Kobs < (31/10)I_240.                     (A28)
\end{verbatim}

The sharper rational endpoints and every nonzero coefficient are
retained in the same records. This is the further original energy input
used in T8 to bound the correction between the actual state-minimum and
energy-minimum sections of the inverse-power family.

\section{The Split-Zero heat comparison transferred to the original
Yang--Mills
Grams}\label{the-split-zero-heat-comparison-transferred-to-the-original-yangmills-grams}

21 September 2026. This is a source-specific transfer, accompanied by a
new calculation of actual Yang--Mills heat correlations. The inspected
Riemann workbench revision is 128aa308dd2a70ba6e073816a957d1ee6414ba2c.
The exact read ranges and source blobs are in
{S\allowbreak{}O\allowbreak{}U\allowbreak{}R\allowbreak{}C\allowbreak{}E\allowbreak{}\_\allowbreak{}I\allowbreak{}N\allowbreak{}T\allowbreak{}A\allowbreak{}K\allowbreak{}E\allowbreak{}.\allowbreak{}j\allowbreak{}s\allowbreak{}o\allowbreak{}n\allowbreak{}.\allowbreak{}}
Its arithmetic quantities remain attached to their own source; the
Yang--Mills quantities below are computed from the original Hamiltonian
and its physical vacuum.

\subsection{T1. What was read and what is
transferred}\label{t1.-what-was-read-and-what-is-transferred}

The complete argument used from the Riemann workbench is HM6-\/-HM14 and
HM19-\/-HM24 of
{H\allowbreak{}E\allowbreak{}A\allowbreak{}T\allowbreak{}\_\allowbreak{}T\allowbreak{}O\allowbreak{}\_\allowbreak{}O\allowbreak{}R\allowbreak{}I\allowbreak{}G\allowbreak{}I\allowbreak{}N\allowbreak{}A\allowbreak{}L\allowbreak{}\_\allowbreak{}Q\allowbreak{}U\allowbreak{}O\allowbreak{}T\allowbreak{}I\allowbreak{}E\allowbreak{}N\allowbreak{}T\allowbreak{}\_\allowbreak{}M\allowbreak{}E\allowbreak{}T\allowbreak{}R\allowbreak{}I\allowbreak{}C\allowbreak{}S\allowbreak{}.\allowbreak{}t\allowbreak{}e\allowbreak{}x\allowbreak{}}
(19 September). It constructs its original coefficient column map Rcal
and a finite heat replacement Rcal\_J=Rcal-T\_s Pcal. The physical
arithmetic source pairing is unchanged and T\_s* T\_s=s\^{}(-1)I. From
the complete original Gram, including both relation cross blocks, HM13
proves the relative column error. Triangle inequalities give
(1-delta)\^{}2 G \textless= G\_J \textless=(1+delta)\^{}2 G. Completing
the square on the same coefficient fibers carries these bounds through
their actual restrictions and quotient minima. HM22-\/-HM24 retain every
determinant rank and give the finite heat depth for a prescribed return
error.

HG6-\/-HG10 of
{O\allowbreak{}R\allowbreak{}I\allowbreak{}G\allowbreak{}I\allowbreak{}N\allowbreak{}A\allowbreak{}L\allowbreak{}\_\allowbreak{}H\allowbreak{}E\allowbreak{}A\allowbreak{}T\allowbreak{}\_\allowbreak{}A\allowbreak{}N\allowbreak{}D\allowbreak{}\_\allowbreak{}S\allowbreak{}O\allowbreak{}U\allowbreak{}R\allowbreak{}C\allowbreak{}E\allowbreak{}\_\allowbreak{}E\allowbreak{}N\allowbreak{}D\allowbreak{}P\allowbreak{}O\allowbreak{}I\allowbreak{}N\allowbreak{}T\allowbreak{}.\allowbreak{}t\allowbreak{}e\allowbreak{}x\allowbreak{}}
were also read. They expand actual
{o\allowbreak{}b\allowbreak{}s\allowbreak{}e\allowbreak{}r\allowbreak{}v\allowbreak{}a\allowbreak{}t\allowbreak{}i\allowbreak{}o\allowbreak{}n\allowbreak{}/\allowbreak{}k\allowbreak{}e\allowbreak{}r\allowbreak{}n\allowbreak{}e\allowbreak{}l\allowbreak{}}
currents with both mixed products and give their error under a relative
change of the same original Gram. The fixed-divisor subexponential
estimate HG1-\/-HG5 is retained as an inspected arithmetic result; no
value of its divisor-dependent constants is assigned to a Yang--Mills
bound.

The common column-map calculation is instantiated below, on actual
Yang--Mills spaces. Its input error is supplied by the original physical
semigroup rather than by the arithmetic dilation. A source-specific map
between arithmetic zeta packets and physical gauge fields is not
asserted. The explicit shared interface and each new map are the
displayed linear maps below.

\subsection{T2. The original physical heat columns and their entire
forcing
defect}\label{t2.-the-original-physical-heat-columns-and-their-entire-forcing-defect}

Fix the original L=2 open box, M=240, a\textgreater0, and
g\^{}2\textgreater=20. Let A=H-E0 on the centered physical Hilbert space
H0, with its unchanged inner product. The source columns R have entries
{r\allowbreak{}p\allowbreak{}=\allowbreak{}(\allowbreak{}W\allowbreak{}p\allowbreak{}-\allowbreak{}\textless{}\allowbreak{}W\allowbreak{}p\allowbreak{}\textgreater{}\allowbreak{}\_\allowbreak{}r\allowbreak{}h\allowbreak{}o\allowbreak{})\allowbreak{}p\allowbreak{}s\allowbreak{}i\allowbreak{}.\allowbreak{}}
From A23,

\begin{verbatim}
Phi=kappa A^(-1)R,
G2=Phi*Phi,   E=Phi*A Phi=kappa G1,
(9/100)I<G2<(13/100)I,
(3kappa/10)I<E<(11kappa/30)I.                           (T1)
\end{verbatim}

All original coefficient axes p are retained. In particular Phi is
injective. The finite-regulator inverse and the original source identity
A Phi=kappa R are established in the analytic proof, not assumed as an
observation metric.

The physical gap needed here has an independent short derivation. In the
original Haar space, minmax applied to B\_xi=K-xi S gives its second
physical eigenvalue \textgreater=3-2Mxi. Its ground eigenvalue is
\textless=0 by the original constant trial vector. Therefore, for the
actual centered A,

\begin{verbatim}
A >=kappa(3-2Mxi)I >=(27kappa/10)I                       (T2)
\end{verbatim}

because xi\textless=1/1600. This is a bound for this original box and
coupling domain. Every volume factor in its derivation is explicit.

For T\textgreater0 define the actual map and its convergent inverse on
H0

\begin{verbatim}
C_T=I-exp(-TA),
C_T^(-1)=sum_(j>=0)exp(-jTA),
||C_T^(-1)|| <=[1-exp(-27kappa T/10)]^(-1).              (T3)
\end{verbatim}

They commute with A on Dom(A); convergence in the graph norm follows by
commuting A through each term and applying the same geometric bound. Put

\begin{verbatim}
Phi_T=C_T Phi=kappa int_0^T exp(-tA)R dt,
E_T=Phi_T*A Phi_T,  G2_T=Phi_T*Phi_T.                   (T4)
\end{verbatim}

Both state and energy errors retain their complete mixed products.
Specifically,

\begin{verbatim}
E_T-E =-Phi*A exp(-TA)Phi
        -Phi*exp(-TA)A Phi
        +Phi*exp(-TA)A exp(-TA)Phi,                     (T5)
\end{verbatim}

with the analogous identity omitting A for G2. The two cross products
are equal here because the original A and its heat semigroup commute;
their origin in the expansion remains explicit. The source equation is

\begin{verbatim}
A Phi_T=kappa R-kappa exp(-TA)R.                        (T6)
\end{verbatim}

Thus the original forcing defect and its primitive are exactly -kappa
exp(-TA)R and -exp(-TA)Phi, respectively.

For epsilon=exp(-27kappa T/10), spectral calculus applied on the
complete centered physical space gives

\begin{verbatim}
(1-epsilon)^2 G2 <= G2_T <= G2,
(1-epsilon)^2 E <= E_T <= E.                            (T7)
\end{verbatim}

For the energy inequality one applies C\_T to A\^{}(1/2)Phi; for the
state inequality one applies it to Phi. These are the actual column-map
instances of HM11, with a sharper one-sided upper bound supplied by the
semigroup. The norm of the omitted column is bounded relatively in its
own original state/energy pairing. All constants are independent of a
after the exact physical time T is inserted; T itself retains its factor
1/kappa.

\subsection{T3. Original support quotients and the corrected
minimum}\label{t3.-original-support-quotients-and-the-corrected-minimum}

For an original face subset F, put V\_F=Phi(C\^{}F), U\_F=A V\_F=kappa
R(C\^{}F). The complex is

\begin{verbatim}
V_F --A--> H0 --0--> 0.                                 (T8)
\end{verbatim}

Its support transition F subset G is the actual inclusion in degree zero
and identity on H0. The differential square commutes by the same
original A. For J=G\textbackslash F, the transported cohomology kernel
is exactly

\begin{verbatim}
V_G/V_F -> ker[H0/U_F -> H0/U_G],
[h] -> [Ah].                                           (T9)
\end{verbatim}

Surjectivity follows by writing a killed class as Ah with h in V\_G.
Changing h by V\_F changes Ah by U\_F, and conversely the injectivity of
A on H0 implies that a change by U\_F has exactly a change by V\_F. The
inverse is {[}Ah{]} -\textgreater{} {[}h{]}. The heat complexes use
V\_F\^{}T=C\_TV\_F with the same target H0 and differential A. The
actual pair (C\_T,C\_T) intertwines their differentials and is
invertible by (T3). It commutes with every support inclusion. Applying
the Split-Zero support reconstruction therefore retains F and its
receiving zero, together with the original killed representative and its
primitive.

The energy-minimum representative of {[}Phi\_J y{]} in V\_G/V\_F has the
old-support coordinate

\begin{verbatim}
x_F=-E_FF^(-1)E_FJ y.
\end{verbatim}

All blocks are taken in the fixed original face coordinates. Its exact
quotient Gram is

\begin{verbatim}
Q_E(F,G)=E_JJ-E_JF E_FF^(-1)E_FJ.                       (T10)
\end{verbatim}

Completing the square in x\_F proves this formula and both inverse
coordinate maps to the quotient. The removed primitive Phi\_F x\_F
remains. The state minimum has the separately calculated formula with
G2; it is not substituted for the energy minimum.

Equation (T7) holds at every vector (x\_F,y) in the same original
coefficient fiber. Taking the infimum in x\_F on both sides proves

\begin{verbatim}
(1-epsilon)^2 Q_E(F,G)<=Q_(E_T)(F,G)<=Q_E(F,G),           (T11)
\end{verbatim}

and the state counterpart. The same proof applies successively to every
finite chain of fixed restrictions and quotient maps, retaining the same
two bounds, not multiplying an error factor at each stage. This is the
actual HM19-\/-HM21 minimum-fiber argument in the physical metric. A
rank-zero quotient retains its unique zero vector and determinant one.

To retain the surrounding physical space, put P\_Phi=Phi G2\^{}(-1)Phi*.
Its multiplication identities give
{P\allowbreak{}\_\allowbreak{}P\allowbreak{}h\allowbreak{}i\allowbreak{}\textasciicircum{}\allowbreak{}2\allowbreak{}=\allowbreak{}P\allowbreak{}\_\allowbreak{}P\allowbreak{}h\allowbreak{}i\allowbreak{}=\allowbreak{}P\allowbreak{}\_\allowbreak{}P\allowbreak{}h\allowbreak{}i\allowbreak{}*\allowbreak{}.\allowbreak{}}
For any original form-domain h perpendicular to im(Phi), the full energy
is

\begin{verbatim}
q_A(Phi c+h)=kappa c*G1 c
            +2kappa Re<Rc,h>+q_A(h).                  (T12)
\end{verbatim}

This is the complete original coupling to the remaining physical
functions. No reducing-subspace property of the finite observation
family is used.

\subsection{T4. A finite heat horizon for every four-return
determinant}\label{t4.-a-finite-heat-horizon-for-every-four-return-determinant}

For a rank-r positive form G and its same-coordinate heat comparison
satisfying (T7), its relative eigenvalues belong to
{[}(1-epsilon)\^{}2,1{]}. The determinant ratio is invariant under the
displayed congruence with G\^{}(-1/2), including its inverse G\^{}(1/2).
Hence

\begin{verbatim}
|log det G_T-log det G| <=2r[-log(1-epsilon)].           (T13)
\end{verbatim}

The difference uses the same original frame and, for energy Grams, the
same kappa\^{}r factor, which cancels exactly in the ratio. For any four
prescribed
{r\allowbreak{}e\allowbreak{}s\allowbreak{}t\allowbreak{}r\allowbreak{}i\allowbreak{}c\allowbreak{}t\allowbreak{}i\allowbreak{}o\allowbreak{}n\allowbreak{}/\allowbreak{}q\allowbreak{}u\allowbreak{}o\allowbreak{}t\allowbreak{}i\allowbreak{}e\allowbreak{}n\allowbreak{}t\allowbreak{}}
returns with coefficients +1,+1,-1,-1 and ranks r\_j\textless=240,

\begin{verbatim}
|L_T-L| <=2 sum_j r_j [-log(1-epsilon)]
          <=1920 epsilon/(1-epsilon).                 (T14)
\end{verbatim}

The final inequality follows by integrating 1/(1-x) from zero to
epsilon. Take the actual physical horizon

\begin{verbatim}
T=10/kappa=5a/g^2.                                     (T15)
\end{verbatim}

Then epsilon=exp(-27). The exact rational exponential enclosure gives

\begin{verbatim}
epsilon<19/10^13,
1920 epsilon/(1-epsilon)<4/10^9.                        (T16)
\end{verbatim}

The calculated inner upper bound is approximately
3.608695327762*10\^{}(-9). The complete rational fraction is in
{g\allowbreak{}e\allowbreak{}n\allowbreak{}e\allowbreak{}r\allowbreak{}a\allowbreak{}t\allowbreak{}e\allowbreak{}d\allowbreak{}/\allowbreak{}h\allowbreak{}e\allowbreak{}a\allowbreak{}t\allowbreak{}\_\allowbreak{}b\allowbreak{}o\allowbreak{}u\allowbreak{}n\allowbreak{}d\allowbreak{}s\allowbreak{}.\allowbreak{}j\allowbreak{}s\allowbreak{}o\allowbreak{}n\allowbreak{}.\allowbreak{}}
This is one common finite heat horizon for all four original returns and
all original support choices in this 240-column family. It requires no
estimate of a newly selected arithmetic comparison constant or
observation angle.

\subsection{T5. Every signed observation current retains its mixed
terms}\label{t5.-every-signed-observation-current-retains-its-mixed-terms}

Fix a surjection Lambda from the original C\^{}240 onto its actual
chosen observation image; for example its rows can read a fixed subset
of original face coordinates. From the energy Gram E define

\begin{verbatim}
Q=(Lambda E^(-1)Lambda*)^(-1),
Omega=Lambda*Q Lambda,
L=E-Omega.                                             (T17)
\end{verbatim}

The minimum representative S=E\^{}(-1)Lambda*Q satisfies Lambda S=I and
E S=Lambda*Q. Direct expansion proves that y*L y is the squared E-norm
of (1-S Lambda)y. Thus both the kernel residual and boundary Gram are
the original canonical ones, with every rectangular mixed entry
preserved. Define the heat versions using E\_T and the identical Lambda.

Put eta=2epsilon-epsilon\^{}2. Equation (T7) gives
(1-eta)E\textless=E\_T\textless=E. Minimize on the identical Lambda
fibers to get the same bounds for Q. For fixed original columns a1,a2,v
and omega\textgreater0 let

\begin{verbatim}
K_i=a_i*L v, B_i=a_i*Omega v,
d_i=a_i*E a_i, E_v=v*E v,
Phi_K=2Re(conj(K1)K2)/omega,
Phi_B=2Re(conj(B1)B2)/omega,
Phi_cross=2Re(conj(K1)B2+conj(B1)K2)/omega.              (T18)
\end{verbatim}

The same-fiber bounds and Cauchy-\/-Schwarz give \textbar Delta
K\_i\textbar\textless=2eta sqrt(d\_i E\_v), \textbar Delta
B\_i\textbar\textless=eta sqrt(d\_i E\_v), while each original
\textbar K\_i\textbar{} and \textbar B\_i\textbar{} is at most sqrt(d\_i
E\_v). Expanding every linear and quadratic product in (T18) proves the
HG8 receiving inequalities

\begin{verbatim}
|Delta Phi_K|<=2sqrt(d1 d2)E_v(4eta+4eta^2)/omega,
|Delta Phi_B|<=2sqrt(d1 d2)E_v(2eta+eta^2)/omega,
|Delta Phi_cross|<=2sqrt(d1 d2)E_v(6eta+4eta^2)/omega.    (T19)
\end{verbatim}

These estimates contain the original phases and both cross products. The
actual source/target numerical values in this application are the
calculated Yang--Mills Grams, and their relative error is (T16). No sign
of a current is inferred from its norm alone.

\subsection{T6. Source and execution
scope}\label{t6.-source-and-execution-scope}

This contribution reads the original HM/HG heat and metric proofs and
proves the displayed Yang--Mills instances in full. The arithmetic
dilation T\_s, its Mellin variables, zeta-zero jets, original source
masses and period-dependent constants remain in the cited RH
construction. Here exp(-TA) is the physical ground-relative semigroup,
with T (T15) in the original time units.

The heat coefficients, original box matrices, graph constants and
rational bounds are independently exercised by the executable tests
described in VERIFICATION.md. Their complete analytic proofs are
supplied as written arguments for review. The tests do not give a Lean
proof, independently audit every older source estimate, or establish the
four-dimensional continuum gap. The standing continuum path and the
growing-volume problem still require bounds beyond the explicitly finite
M-dependent domain proved here.

\subsection{T7. The latest inverse-power observation has a concrete
receiving
instance}\label{t7.-the-latest-inverse-power-observation-has-a-concrete-receiving-instance}

After the heat-source intake, the complete 21 September source
{I\allowbreak{}N\allowbreak{}V\allowbreak{}E\allowbreak{}R\allowbreak{}S\allowbreak{}E\allowbreak{}\_\allowbreak{}P\allowbreak{}O\allowbreak{}W\allowbreak{}E\allowbreak{}R\allowbreak{}\_\allowbreak{}C\allowbreak{}O\allowbreak{}N\allowbreak{}D\allowbreak{}U\allowbreak{}C\allowbreak{}T\allowbreak{}O\allowbreak{}R\allowbreak{}\_\allowbreak{}O\allowbreak{}B\allowbreak{}S\allowbreak{}E\allowbreak{}R\allowbreak{}V\allowbreak{}A\allowbreak{}B\allowbreak{}I\allowbreak{}L\allowbreak{}I\allowbreak{}T\allowbreak{}Y\allowbreak{}.\allowbreak{}t\allowbreak{}e\allowbreak{}x\allowbreak{}}
was read at the same pinned commit. Its IK9-\/-IK10 construct a literal
left inverse of selected original conductor rows on their inverse-power
space. IK12 returns that inverse through the original source and
observation Grams, retaining their eigenvalue factors. IK13 explicitly
permits its lower bound to tend to zero with its cutoff. Those
coefficients, root differences and guards remain in that construction.

In the present Yang--Mills family the actual inverse-power source is

\begin{verbatim}
Phi=kappa A^(-1)R:C^240->H0,
O=R*:H0->C^240,
O Phi=G^(1).                                           (T20)
\end{verbatim}

The independent computed bound (A23) makes this original response matrix
invertible, so its exact coefficient left inverse is

\begin{verbatim}
W=(G^(1))^(-1) R*,
W Phi=I_240,   ||(G^(1))^(-1)||<=10/3.                  (T21)
\end{verbatim}

This is a proved inverse on these specific columns, not an assumed
observation angle. The original quotient metric for the observation O is
(G\^{}(0))\^{}(-1): minimizing the physical Hilbert norm over R*h=z has
the representative R(G\^{}(0))\^{}(-1)z. Its full residual is
perpendicular to im(R), as multiplication by R* verifies. Thus
{P\allowbreak{}\_\allowbreak{}R\allowbreak{}=\allowbreak{}R\allowbreak{}(\allowbreak{}G\allowbreak{}\textasciicircum{}\allowbreak{}(\allowbreak{}0\allowbreak{})\allowbreak{})\allowbreak{}\textasciicircum{}\allowbreak{}(\allowbreak{}-\allowbreak{}1\allowbreak{})\allowbreak{}R\allowbreak{}*\allowbreak{}}
is the actual orthogonal projection, and every nonzero coefficient
vector c satisfies

\begin{verbatim}
||P_R Phi c||^2=c*G^(1)(G^(0))^(-1)G^(1)c,
||Phi c||^2=c*G^(2)c.                                  (T22)
\end{verbatim}

Returning the three computed original metrics (A23),(A25) gives

\begin{verbatim}
G^(1)(G^(0))^(-1)G^(1) >= (3/34)I,
||P_R Phi c||^2 / ||Phi c||^2 > 150/221.                 (T23)
\end{verbatim}

Indeed the first coefficient is
{(\allowbreak{}3\allowbreak{}/\allowbreak{}1\allowbreak{}0\allowbreak{})\allowbreak{}\textasciicircum{}\allowbreak{}2\allowbreak{}/\allowbreak{}(\allowbreak{}5\allowbreak{}1\allowbreak{}/\allowbreak{}5\allowbreak{}0\allowbreak{})\allowbreak{}=\allowbreak{}3\allowbreak{}/\allowbreak{}3\allowbreak{}4\allowbreak{},\allowbreak{}}
and dividing by the upper state coefficient 13/100 gives 150/221. Every
source metric and its inverse has been retained. This is the IK12
{l\allowbreak{}e\allowbreak{}f\allowbreak{}t\allowbreak{}-\allowbreak{}i\allowbreak{}n\allowbreak{}v\allowbreak{}e\allowbreak{}r\allowbreak{}s\allowbreak{}e\allowbreak{}/\allowbreak{}m\allowbreak{}i\allowbreak{}n\allowbreak{}i\allowbreak{}m\allowbreak{}u\allowbreak{}m\allowbreak{}}
argument applied to the original physical inverse-power family, with the
needed Gram bounds actually calculated here. It detects all 240 columns
simultaneously at L=2, every a\textgreater0 and g\^{}2\textgreater=20.
It does not supply an infinite-volume fraction or an independence
assertion for an uncomputed concatenation of higher powers.

The portion outside this observation is the exact vector (I-P\_R)Phi c.
Its squared norm is

\begin{verbatim}
c*[G^(2)-G^(1)(G^(0))^(-1)G^(1)]c,                       (T24)
\end{verbatim}

and is positive semidefinite by the original orthogonal decomposition.
No component is removed from the energy formula (T12). The maps
(T20)-\/-(T24) state the entire receiving construction; they make no
identification of arithmetic zeta zeros with physical spectral points.

\subsection{T8. The actual residual between the state and energy
sections is
bounded}\label{t8.-the-actual-residual-between-the-state-and-energy-sections-is-bounded}

Retain the original orthogonal state projection P\_R from T7 and define

\begin{verbatim}
h=(I-P_R)Phi,
B0=(G^(0))^(-1)G^(1),   P_R Phi=R B0.                  (T25)
\end{verbatim}

All columns of h satisfy R*h=0. Since A Phi=kappa R, the mixed energy
q\_A(Phi c,h d) vanishes for every coefficient pair c,d, by the original
Hilbert adjoint identity. Thus Phi is the actual minimum-energy section
of the affine fibers R*f=G\^{}(1)c. The state-minimum section of that
same fiber is R B0 c. They have the exact comparison

\begin{verbatim}
h*Ah=kappa[G^(1)(G^(0))^(-1)Kobs(G^(0))^(-1)G^(1)-G^(1)],
(P_R Phi)*A(P_R Phi)=Phi*A Phi+h*Ah.                    (T26)
\end{verbatim}

To verify the first identity, expand the whole expression
(Phi-RB0)*A(Phi-RB0). Both cross products are kappa G\^{}(1), while the
last term is kappa B0* Kobs B0; their signed sum is (T26). This also
proves nonnegativity of the right side. The state residual remains the
matrix (T24).

The three raw bounds (A23),(A25),(A28) give the strict original energy
bound

\begin{verbatim}
0<=h*Ah < (1322/7203) Phi*A Phi.                        (T27)
\end{verbatim}

Indeed congruence by (G\^{}(1))\^{}(-1/2), with its explicitly inverse
square root, puts the bracket in (T26) below

\begin{verbatim}
[(31/10)(50/49)^2(11/30)-1]I = (1322/7203)I.
\end{verbatim}

Returning that congruence proves (T27) in the original coefficient
coordinates. Thus the state-section correction has at most this fraction
(less than0.184) of the original primitive\textquotesingle s energy, and
(T23) retains more than150/221 of its state norm in the actual observed
force span. The difference itself remains an element of ker(R*) with its
complete state and energy Grams, rather than being dropped. Both
conclusions hold simultaneously for all240 original columns on the
specified finite box and coupling domain.

For a prescribed observation z, the two exact sections are
Phi(G\^{}(1))\^{}(-1)z and R(G\^{}(0))\^{}(-1)z, and their difference is
h(G\^{}(1))\^{}(-1)z. Their source-to-observation compositions are both
the identity. This supplies the explicit section-correction morphism and
its quantitative energy return in the original metric. It is the same
minimum-section mechanism read in the Split-Zero source, now with all
physical Grams evaluated and the entire correction bounded on its stated
domain.

\section{Original Yang--Mills heat: a box-independent analytic radius
and row-sum
remainder}\label{original-yangmills-heat-a-box-independent-analytic-radius-and-row-sum-remainder}

21 September 2026. This continuation addresses the two quantities left
by the heat edition: its radius 1/M and the coupling of its observed
states to the full physical complement. All operator coefficients, Haar
masses, original link and plaquette coordinates, and state and energy
pairings are retained. This note proves the local-source-to-heat
argument used here in full. It does not certify every earlier
stronger-coupling assertion. Its coefficient tables are the inherited,
separately replayed heat tables; new constant and matrix calculations
have their own records. No continuum mass-gap conclusion,
historical-priority determination, or external analytical review is
asserted.

\subsection{U1. Original operator, physical time, and coefficient
maps}\label{u1.-original-operator-physical-time-and-coefficient-maps}

For each L\textgreater=2 take vertices \{-L,...,L\}\^{}3, all contained
positive nearest-neighbor links e=(n,i), and all elementary faces
p=(n;i,j), i\textless j. Retain

\begin{verbatim}
Wp=tr(U_i(n) U_j(n+e_i) U_i(n+e_j)^(-1) U_j(n)^(-1)),
T_a=-i sigma_a/2,  X_e,a f=d/dt f(...,exp(t T_a)U_e,...)|0,
K=-sum_(e,a) X_e,a^2,
H(x)=kappa K+kappa sum_p x_p(2-Wp),
kappa=2g^2/a,   x_p=xi=1/(4g^4),   a,g>0.                 (U1)
\end{verbatim}

The physical Hilbert space is the invariant subspace of original product
Haar probability L2(dU), with every vertex including boundary vertices
averaged. H has physical H2 domain and H1 form domain, since the
original potential is smooth bounded on this finite compact group. Its
smooth positive unit ground vector psi and simple bottom energy E0
follow from the modulus form inequality, elliptic regularity, the
maximum principle, and the identity

\begin{verbatim}
q_(H-E0)(psi f)=kappa int psi^2 sum_i |X_i f|^2.         (U2)
\end{verbatim}

The multiplication map f -\textgreater{} psi f and inverse division by
psi are unitary between L2(rho dU) and original L2(dU), rho=psi\^{}2.
The time coordinates are tau=kappa t and t=tau/kappa. No field variable
is rescaled.

For a product-spin label j=(j\_e), define pi\_j(U)=tensor\_e
pi\_(j\_e)(U\_e),
{d\allowbreak{}\_\allowbreak{}j\allowbreak{}=\allowbreak{}p\allowbreak{}r\allowbreak{}o\allowbreak{}d\allowbreak{}u\allowbreak{}c\allowbreak{}t\allowbreak{}\_\allowbreak{}e\allowbreak{}(\allowbreak{}2\allowbreak{}j\allowbreak{}\_\allowbreak{}e\allowbreak{}+\allowbreak{}1\allowbreak{})\allowbreak{},\allowbreak{}}
c\_j=sum\_e j\_e(j\_e+1), and the actual coefficient

\begin{verbatim}
A_j=d_j int f(U) pi_j(U)^* dU,
f(U)=sum_j Tr(A_j pi_j(U)).                             (U3)
\end{verbatim}

This Fourier map and its inverse follow from the original Haar
matrix-coefficient orthogonality. The trace-coefficient norm is
\textbar\textbar f\textbar\textbar\_X=sum\_j
\textbar\textbar A\_j\textbar\textbar\_1. Let X0 be its zero-Haar
physical part, and Y0 its dense domain
{s\allowbreak{}u\allowbreak{}m\allowbreak{}\_\allowbreak{}(\allowbreak{}j\allowbreak{}!\allowbreak{}=\allowbreak{}0\allowbreak{})\allowbreak{}c\allowbreak{}\_\allowbreak{}j\allowbreak{}|\allowbreak{}|\allowbreak{}A\allowbreak{}\_\allowbreak{}j\allowbreak{}|\allowbreak{}|\allowbreak{}\_\allowbreak{}1\allowbreak{}\textless{}\allowbreak{}i\allowbreak{}n\allowbreak{}f\allowbreak{}i\allowbreak{}n\allowbreak{}i\allowbreak{}t\allowbreak{}y\allowbreak{}.\allowbreak{}}
The actual K maps Y0 bijectively onto X0 by multiplying each coefficient
by c\_j; its inverse divides it by that same c\_j. These are auxiliary
convergence estimates on the original coefficients, not replacements for
(U2).

The compact-group Fourier algebra and its product norm are classical
(Eymard, 1964). Here their required inequalities are proved directly. In
a product of two coefficient functions, decompose pi\_j tensor pi\_k
into its full irreducible and multiplicity spaces by a unitary
intertwiner. Pinching to all its diagonal representation blocks and
tracing each multiplicity space returns the exact output coefficients.
Pinching is an average of unitary conjugations. For a partial trace,
trace-norm duality gives \textbar Tr{[}(Z tensor
I)C{]}\textbar\textless=\textbar\textbar C\textbar\textbar\_1 when
\textbar\textbar Z\textbar\textbar\textless=1. Hence

\begin{verbatim}
||fh||_X<=||f||_X||h||_X,
||A(X_e,a f_j)||_1<=j_e||A_j||_1.                      (U4)
\end{verbatim}

Every representation dimension and every multiplicity is included in
this argument.

\subsection{U2. The graph bound that enters the source
estimate}\label{u2.-the-graph-bound-that-enters-the-source-estimate}

A nonzero gauge-invariant coefficient has an invariant tensor on the
original incident indices at every vertex. Grouping the factors gives an
intertwiner from the dual spin-j\_e factor into the tensor product of
the others. The maximum target weight is their summed spin, so

\begin{verbatim}
j_e<=sum_(f incident v, f!=e) j_f.                      (U5)
\end{verbatim}

Fix e=\{u,v\}, and let E1 be the other edges at its endpoints. Because
the graph is simple and triangle-free, their outer endpoints are
distinct. Put J1=sum\_(f in E1)j\_f. The two endpoint inequalities give
J1\textgreater=2j\_e. Sum the outer-endpoint inequalities: every edge
outside E1 and \{e\} is counted at most twice, so J1\textless=2J2, where
J2 counts these outer edges once. Therefore sum\_f
{j\allowbreak{}\_\allowbreak{}f\allowbreak{}\textgreater{}\allowbreak{}=\allowbreak{}j\allowbreak{}\_\allowbreak{}e\allowbreak{}+\allowbreak{}J\allowbreak{}1\allowbreak{}+\allowbreak{}J\allowbreak{}2\allowbreak{}\textgreater{}\allowbreak{}=\allowbreak{}4\allowbreak{}j\allowbreak{}\_\allowbreak{}e\allowbreak{}.\allowbreak{}}
The original half-integer spectrum obeys j(j+1)\textgreater=(3/2)j, and
consequently

\begin{verbatim}
c_j>=6j_e on every physical block;
sum_e j_e<=(2/3)c_j on every product block.             (U6)
\end{verbatim}

In particular c\_j\textgreater=3 on each nonconstant physical block. The
fundamental elementary plaquette has c\_j=3. This argument includes
active bridges and boundary vertices; no cycle decomposition of every
active edge is presumed.

\subsection{U3. Local labelled source and its exact assembly
kernel}\label{u3.-local-labelled-source-and-its-exact-assembly-kernel}

For every nonempty finite edge label S retain a zero-Haar physical
coefficient family A\_(S,j), with supp(j) contained in S; the label may
exceed the active support. Put

\begin{verbatim}
||v||loc=sup_e sum_(S containing e,j!=0)c_j||A_(S,j)||_1. (U7)
\end{verbatim}

On a finite graph this is a complete space, and its total c-weighted
coefficient sum is at most \textbar E\_L\textbar{}
\textbar\textbar v\textbar\textbar loc. Assembly sends A\_(S,j) to
sum\_(S containing supp(j)) A\_(S,j). Its kernel is exactly the family
equations that all those sums vanish. For each S strictly containing
supp(j), the relation with coefficient A at (S,j) and -A at (supp(j),j)
has zero assembly. The actual coefficients at all enlarged S give a
unique decomposition of every assembly-kernel element into these
relations. Its inverse reads those same enlarged-label coefficients.
This supplies the support-preserving quotient map and both inverse laws.

Define, retaining the original union label,

\begin{verbatim}
B(v,h)_S=K^(-1) Q_H sum_(S1 union S2=S) Gamma(v_S1,h_S2),
Gamma(f,h)=sum_(e,a)(X_e,a f)(X_e,a h),  Q_H=I-P_H.     (U8)
\end{verbatim}

Each Haar scalar removed by Q\_H is recorded separately. For an anchor
in S1, (U4) and (U6) bound the output before K inversion by

\begin{verbatim}
3 sum_(S1 containing anchor,j) ||A_(S1,j)||_1
    sum_e j_e sum_(S2 containing e,k) k_e||B_(S2,k)||_1
  <=(1/3)||v||loc||h||loc.
\end{verbatim}

The anchor-in-S2 contribution gives the other 1/3. The original output
Casimir cancels precisely against its nonzero inverse Casimir in (U8).
Thus

\begin{verbatim}
||B(v,h)||loc<=(2/3)||v||loc||h||loc.                   (U9)
\end{verbatim}

\subsection{U4. Exact elementary coefficient and convergent complex
source}\label{u4.-exact-elementary-coefficient-and-convergent-complex-source}

For the original four-link word, the coefficient matrix on
(C2)\^{}tensor4 has entries

\begin{verbatim}
A_(i1,i2,1-i2,1-i3 ; i0,i1,1-i3,1-i0) += (-1)^(i0+i2),
i0,i1,i2,i3 in {0,1}.                                  (U10)
\end{verbatim}

The identity
{(\allowbreak{}U\allowbreak{}\textasciicircum{}\allowbreak{}(\allowbreak{}-\allowbreak{}1\allowbreak{})\allowbreak{})\allowbreak{}\_\allowbreak{}(\allowbreak{}r\allowbreak{},\allowbreak{}c\allowbreak{})\allowbreak{}=\allowbreak{}(\allowbreak{}-\allowbreak{}1\allowbreak{})\allowbreak{}\textasciicircum{}\allowbreak{}(\allowbreak{}r\allowbreak{}+\allowbreak{}c\allowbreak{})\allowbreak{}}
U\_(1-c,1-r) proves its trace formula. Its four nonzero 2-by-2 blocks,
up to row/column order, are {[}{[}1,-1{]},{[}-1,1{]}{]}. Direct
multiplication gives (A* A)\^{}2=4A* A and Tr(A* A)=16. Thus
\textbar A\textbar=A* A/2 and \textbar\textbar Wp\textbar\textbar\_X=8.
At most four original plaquettes meet an edge.

Allow independent complex x\_p with
max\_p\textbar x\_p\textbar\textless=r. The source equation is

\begin{verbatim}
v=(1/3)sum_p x_p Wp+B(v,v).                            (U11)
\end{verbatim}

Its homogeneous order-n coefficient obeys the original binary
recurrence, giving

\begin{verbatim}
a_n(r)=Cat_(n-1)(2/3)^(n-1)(32r)^n,
sum_(n>=1) a_n(r)=r_v(r)=(3/4)(1-sqrt(1-256r/3)).       (U12)
\end{verbatim}

The radius used throughout the new bound is exactly

\begin{verbatim}
R0=3/256.                                              (U13)
\end{verbatim}

For r\textless R0 the series is norm-analytic. It converges absolutely
also at r=R0, where its tail is
{(\allowbreak{}3\allowbreak{}/\allowbreak{}4\allowbreak{})\allowbreak{}b\allowbreak{}i\allowbreak{}n\allowbreak{}o\allowbreak{}m\allowbreak{}(\allowbreak{}2\allowbreak{}N\allowbreak{},\allowbreak{}N\allowbreak{})\allowbreak{}/\allowbreak{}4\allowbreak{}\textasciicircum{}\allowbreak{}N\allowbreak{}\textless{}\allowbreak{}=\allowbreak{}3\allowbreak{}/\allowbreak{}(\allowbreak{}4\allowbreak{}s\allowbreak{}q\allowbreak{}r\allowbreak{}t\allowbreak{}(\allowbreak{}N\allowbreak{}+\allowbreak{}1\allowbreak{})\allowbreak{})\allowbreak{}.\allowbreak{}}
The finite coefficient proof is
{C\allowbreak{}a\allowbreak{}t\allowbreak{}\_\allowbreak{}N\allowbreak{}/\allowbreak{}4\allowbreak{}\textasciicircum{}\allowbreak{}N\allowbreak{}=\allowbreak{}2\allowbreak{}(\allowbreak{}b\allowbreak{}\_\allowbreak{}N\allowbreak{}-\allowbreak{}b\allowbreak{}\_\allowbreak{}(\allowbreak{}N\allowbreak{}+\allowbreak{}1\allowbreak{})\allowbreak{})\allowbreak{},\allowbreak{}}
{b\allowbreak{}\_\allowbreak{}N\allowbreak{}=\allowbreak{}b\allowbreak{}i\allowbreak{}n\allowbreak{}o\allowbreak{}m\allowbreak{}(\allowbreak{}2\allowbreak{}N\allowbreak{},\allowbreak{}N\allowbreak{})\allowbreak{}/\allowbreak{}4\allowbreak{}\textasciicircum{}\allowbreak{}N\allowbreak{}.\allowbreak{}}
No strict contraction at the boundary is used.

Each order-n label is an edge-connected union of at most n plaquettes.
The coefficients agree literally in every original box containing that
union. This follows inductively from (U8): derivatives require a shared
original edge, and K inversion acts on the same coefficients.

On a finite graph the norm convergence controls every first and second
coordinate derivative, since j\_e j\_f\textless=c\_j. Thus the actual
assembled function solves

\begin{verbatim}
Kv=sum_p x_pWp+Gamma(v,v)-C(x),
C(x)=P_H Gamma(v,v).                                   (U14)
\end{verbatim}

For real x,
{p\allowbreak{}s\allowbreak{}i\allowbreak{}=\allowbreak{}e\allowbreak{}x\allowbreak{}p\allowbreak{}(\allowbreak{}v\allowbreak{})\allowbreak{}/\allowbreak{}s\allowbreak{}q\allowbreak{}r\allowbreak{}t\allowbreak{}(\allowbreak{}i\allowbreak{}n\allowbreak{}t\allowbreak{}}
exp(2v)) is positive and solves the original eigen-equation with

\begin{verbatim}
E0=kappa[2sum_p x_p-C(x)].                              (U15)
\end{verbatim}

Equation (U2) then proves it is the actual ground vector and fixes every
scalar. Elliptic bootstrapping makes it smooth. For complex x the
nonzero function exp(v) is still a smooth solution of that
eigen-equation. Its squared integral is addressed below, rather than
presumed nonzero.

\subsection{U5. The full drift, including its scalar
row}\label{u5.-the-full-drift-including-its-scalar-row}

Set D\_v f=2 Gamma(v,f). Equations (U4),(U6) give, on the original
coefficient spaces,

\begin{verbatim}
||D_v h||_X<=chi(r)||Kh||_X,
chi(r)=(2/3)r_v(r)=(1-sqrt(1-256r/3))/2<=1/2.           (U16)
\end{verbatim}

Both the constant and nonconstant outputs are included. Indeed the sum
over generators is at most 6 sum\_j
\textbar\textbar A\_j(h)\textbar\textbar\_1 sum\_e j\_e {[}r\_v/6{]},
and
{s\allowbreak{}u\allowbreak{}m\allowbreak{}\_\allowbreak{}e\allowbreak{}j\allowbreak{}\_\allowbreak{}e\allowbreak{}\textless{}\allowbreak{}=\allowbreak{}(\allowbreak{}2\allowbreak{}/\allowbreak{}3\allowbreak{})\allowbreak{}c\allowbreak{}\_\allowbreak{}j\allowbreak{}.\allowbreak{}}
Define

\begin{verbatim}
T=Q_H D_v K^(-1):X0->X0,
p=P_H D_v K^(-1):X0->C,
S=(I-T)^(-1)=sum_(n>=0)T^n.                            (U17)
\end{verbatim}

The stacked bound is
{|\allowbreak{}p\allowbreak{}y\allowbreak{}|\allowbreak{}+\allowbreak{}|\allowbreak{}|\allowbreak{}T\allowbreak{}y\allowbreak{}|\allowbreak{}|\allowbreak{}\_\allowbreak{}X\allowbreak{}\textless{}\allowbreak{}=\allowbreak{}c\allowbreak{}h\allowbreak{}i\allowbreak{}|\allowbreak{}|\allowbreak{}y\allowbreak{}|\allowbreak{}|\allowbreak{}\_\allowbreak{}X\allowbreak{};\allowbreak{}}
in particular \textbar\textbar S\textbar\textbar\textless=1/(1-chi).
Define the actual mean functional by

\begin{verbatim}
mu_x(F)=P_HF+p S Q_HF.                                 (U18)
\end{verbatim}

For chi\textless=1/2,
{|\allowbreak{}m\allowbreak{}u\allowbreak{}\_\allowbreak{}x\allowbreak{}(\allowbreak{}F\allowbreak{})\allowbreak{}|\allowbreak{}\textless{}\allowbreak{}=\allowbreak{}|\allowbreak{}|\allowbreak{}F\allowbreak{}|\allowbreak{}|\allowbreak{}\_\allowbreak{}X\allowbreak{}}
and
{|\allowbreak{}m\allowbreak{}u\allowbreak{}\_\allowbreak{}x\allowbreak{}(\allowbreak{}Q\allowbreak{}\_\allowbreak{}H\allowbreak{}F\allowbreak{})\allowbreak{}|\allowbreak{}\textless{}\allowbreak{}=\allowbreak{}c\allowbreak{}h\allowbreak{}i\allowbreak{}|\allowbreak{}|\allowbreak{}Q\allowbreak{}\_\allowbreak{}H\allowbreak{}F\allowbreak{}|\allowbreak{}|\allowbreak{}\_\allowbreak{}X\allowbreak{}/\allowbreak{}(\allowbreak{}1\allowbreak{}-\allowbreak{}c\allowbreak{}h\allowbreak{}i\allowbreak{})\allowbreak{}.\allowbreak{}}
The associated source inverse is h=K\^{}(-1)S Q\_HF in Y0, with

\begin{verbatim}
(K-D_v)h=F-mu_x(F).                                    (U19)
\end{verbatim}

All maps are analytic for r\textless R0.

Here (U18) equals the actual complex density functional as well. Put
Z\_x=int exp(2v). Integration by parts gives int exp(2v)(K-D\_v)h=0 for
h in Y0, since its first and second derivatives converge. Equation (U19)
therefore gives

\begin{verbatim}
int exp(2v)F=mu_x(F) Z_x.                               (U20)
\end{verbatim}

If Z\_x vanished, the left side would vanish for every physical Fourier
polynomial. Their uniform closure contains the conjugate of the nonzero
smooth gauge-invariant exp(2v). Its squared modulus would then have zero
integral, a contradiction. Hence Z\_x is nonzero, and mu\_x(F)=int
exp(2v)F/Z\_x. For real x it is precisely expectation in rho. The exact
bilinear integration identity is

\begin{verbatim}
mu_x((K-D_v)h . G)=mu_x(Gamma(h,G)).                    (U21)
\end{verbatim}

Complex x uses the bilinear analytic continuation; the physical
real-coupling Hilbert pairing retains conjugation in its first entry.

\subsection{U6. Closed generator and heat flow in the same
coefficients}\label{u6.-closed-generator-and-heat-flow-in-the-same-coefficients}

Let L\_x=K-Q\_H D\_v on X0, Dom(L\_x)=Y0. The graph norms of K and L\_x
are equivalent because \textbar\textbar Q\_H
{D\allowbreak{}\_\allowbreak{}v\allowbreak{}h\allowbreak{}|\allowbreak{}|\allowbreak{}\textless{}\allowbreak{}=\allowbreak{}c\allowbreak{}h\allowbreak{}i\allowbreak{}|\allowbreak{}|\allowbreak{}K\allowbreak{}h\allowbreak{}|\allowbreak{}|\allowbreak{},\allowbreak{}}
chi\textless1. Thus L\_x is closed and densely defined. For
h\textgreater0,

\begin{verbatim}
I+hL_x=[I-hQ_H D_v(I+hK)^(-1)](I+hK),                 (U22)
\end{verbatim}

and the square bracket is invertible by a Neumann series of norm ratio
at most chi. The original sum of trace norms also gives

\begin{verbatim}
||(I+hL_x)f||_X
  >=||f||_X+h(1-chi)||Kf||_X
  >=[1+3h(1-chi)]||f||_X.                              (U23)
\end{verbatim}

All powers of the actual resolvent obey the corresponding product bound.
The
{H\allowbreak{}i\allowbreak{}l\allowbreak{}l\allowbreak{}e\allowbreak{}–\allowbreak{}Y\allowbreak{}o\allowbreak{}s\allowbreak{}i\allowbreak{}d\allowbreak{}a\allowbreak{}/\allowbreak{}L\allowbreak{}u\allowbreak{}m\allowbreak{}e\allowbreak{}r\allowbreak{}–\allowbreak{}P\allowbreak{}h\allowbreak{}i\allowbreak{}l\allowbreak{}l\allowbreak{}i\allowbreak{}p\allowbreak{}s\allowbreak{}}
generation theorem applies to this explicitly closed dense domain and
the proved surjectivity (U22). It gives

\begin{verbatim}
E_x(tau)=exp(-tau L_x),
||E_x(tau)||<=exp[-d(r)tau],  d(r)=3(1-chi(r))>=3/2.    (U24)
\end{verbatim}

For completeness, the shifted operator L\_x-d(r) is dissipative with the
negative-generator convention: the directional norm derivative of -Kf is
-\textbar\textbar Kf\textbar\textbar, and its perturbation is at most
chi\textbar\textbar Kf\textbar\textbar. Its sufficiently large positive
resolvent is obtained by (U22) and the resolvent identity. This proves
the decay constant, rather than using a norm for one inverse as a
substitute for a semigroup estimate. Lumer--Phillips (1961), Theorem3.1,
pp686--687, is the primary generation reference.

Each resolvent in (U22) is analytic in x. Its exponential Euler powers
are uniformly bounded on compact complex source polydisks and converge
strongly to (U24). Here the latter assertion has a direct integral
proof:
{(\allowbreak{}I\allowbreak{}+\allowbreak{}h\allowbreak{}L\allowbreak{})\allowbreak{}\textasciicircum{}\allowbreak{}(\allowbreak{}-\allowbreak{}1\allowbreak{})\allowbreak{}f\allowbreak{}=\allowbreak{}i\allowbreak{}n\allowbreak{}t\allowbreak{}\_\allowbreak{}0\allowbreak{}\textasciicircum{}\allowbreak{}i\allowbreak{}n\allowbreak{}f\allowbreak{}i\allowbreak{}n\allowbreak{}i\allowbreak{}t\allowbreak{}y\allowbreak{}}
exp(-u)E\_x(hu)f du, by integration on the original generator domain and
density. The m-fold product at h=tau/m is the average of E\_x(tau
Y\_m/m)f, where Y\_m has the convolution density
{e\allowbreak{}x\allowbreak{}p\allowbreak{}(\allowbreak{}-\allowbreak{}u\allowbreak{})\allowbreak{}u\allowbreak{}\textasciicircum{}\allowbreak{}(\allowbreak{}m\allowbreak{}-\allowbreak{}1\allowbreak{})\allowbreak{}/\allowbreak{}(\allowbreak{}m\allowbreak{}-\allowbreak{}1\allowbreak{})\allowbreak{}!\allowbreak{},\allowbreak{}}
mean m and variance m. The nonnegative scalar density has mass one.
Chebyshev\textquotesingle s inequality outside a fixed neighborhood of
tau, strong continuity inside it, and the contraction bound give the
stated strong convergence. No spectral diagonalization of the perturbed
coefficient generator is presumed. Applying scalar Vitali convergence to
every bounded functional and then the locally bounded weak-to-strong
holomorphy principle proves analytic dependence of E\_x(tau).
Equivalently the sectorial resolvent obtained from (U16) gives the same
contour construction. This concerns source analyticity, not a
holomorphic continuation of physical time across arbitrary rays.

For real x the full physical heat return is

\begin{verbatim}
T_x(tau)F=E_x(tau)Q_HF+mu_x(F)-mu_x(E_x(tau)Q_HF).      (U25)
\end{verbatim}

On the dense smooth domain it solves the actual equation with generator
K-D\_v, preserves its original mean, and agrees with the physical
Hilbert semigroup by uniqueness. Its positive physical time is
tau/kappa. The same formula defines the analytic continuation for
complex x. In particular

\begin{verbatim}
||T_x(tau)F-mu_x(F)||_X
  <=exp[-d(r)tau]||Q_HF||_X/(1-chi).                   (U26)
\end{verbatim}

Every physical eigenfunction in a finite box is smooth and belongs to Y0
after subtracting its Haar constant. This last membership follows from
Plancherel and Cauchy--Schwarz with a sufficiently high original Casimir
power: sum\_j
{d\allowbreak{}\_\allowbreak{}j\allowbreak{}\textasciicircum{}\allowbreak{}2\allowbreak{}(\allowbreak{}1\allowbreak{}+\allowbreak{}c\allowbreak{}\_\allowbreak{}j\allowbreak{})\allowbreak{}\textasciicircum{}\allowbreak{}(\allowbreak{}-\allowbreak{}s\allowbreak{})\allowbreak{}\textless{}\allowbreak{}i\allowbreak{}n\allowbreak{}f\allowbreak{}i\allowbreak{}n\allowbreak{}i\allowbreak{}t\allowbreak{}y\allowbreak{}}
for s\textgreater3\textbar E\_L\textbar/2. Thus (U24) applied to the
exact eigen-equation proves A=H-E0\textgreater=kappa d(r) on the
complete centered physical Hilbert space. Compact spectral resolution
gives the full form-domain statement. This is a rederivation of the
elementary strong-coupling lower bound on this explicit domain, not of
every stronger historical source estimate.

\subsection{U7. The row-sum estimate that removes the volume
factor}\label{u7.-the-row-sum-estimate-that-removes-the-volume-factor}

For fixed original p put q\_tau=E\_x(tau)Wp and h\_tau=K\^{}(-1)S
q\_tau. The exact mean equation gives

\begin{verbatim}
(K-D_v)h_tau=T_x(tau)Wp-mu_x(Wp).                      (U27)
\end{verbatim}

For arbitrary complex coefficients z\_q with
max\_q\textbar z\_q\textbar\textless=1, set G\_z=sum\_q z\_qWq. The
original edge incidences and trace coefficient 8 give

\begin{verbatim}
sup_e sum_(q containing e) (1/2)||z_q Wq||_X<=16.
\end{verbatim}

Consequently, with every generator retained,

\begin{verbatim}
||Gamma(h,G_z)||_X
  <=3*16 sum_j sum_e j_e||A_j(h)||_1
  <=32||Kh||_X.                                       (U28)
\end{verbatim}

Use (U21), \textbar\textbar mu\_x\textbar\textbar\textless=1,
\textbar\textbar S\textbar\textbar\textless=1/(1-chi), and
\textbar\textbar Wp\textbar\textbar\_X=8. The connected analytic heat
matrix satisfies

\begin{verbatim}
Chat_pq(tau;x)=mu_x[(T_x(tau)Wp-mu_xWp)Wq],
sum_q |Chat_pq(tau;x)|
  <=[256/(1-chi(r))] exp[-d(r)tau]
  <=512 exp(-3tau/2).                                  (U29)
\end{verbatim}

The sum equality to a supremum over phases z\_q is exact in every
original finite face set. At real x this is precisely
\textless r\_p,exp(-tau A/kappa)r\_q\textgreater. Its transpose symmetry
extends to complex x by analyticity from the real multivariate polydisk.
Thus the same bound holds for columns and for the coefficient-space
operator norm. No M is introduced.

There is a useful sharper bound that uses the original Haar projection
before estimating the returned mean. The physical Fourier block with
spin one-half on the four edges of a specified elementary plaquette and
zero on all other edges is one-dimensional: at each two-edge vertex the
invariant contraction is unique, and following the original indices
gives Wp. Consequently, for physical h in Y0,

\begin{verbatim}
|P_H Gamma(h,G_z)| <= (1/8)||Kh||_X.
\end{verbatim}

Indeed the contributing block is a\_p Wp, has
norm8\textbar a\_p\textbar{} and Casimir3, and its Haar pairing with Wp
is a\_p. All other product-spin labels integrate to zero. On the
nonconstant output, (U18) gives
{|\allowbreak{}(\allowbreak{}m\allowbreak{}u\allowbreak{}-\allowbreak{}P\allowbreak{}\_\allowbreak{}H\allowbreak{})\allowbreak{}Q\allowbreak{}F\allowbreak{}|\allowbreak{}\textless{}\allowbreak{}=\allowbreak{}c\allowbreak{}h\allowbreak{}i\allowbreak{}|\allowbreak{}|\allowbreak{}Q\allowbreak{}F\allowbreak{}|\allowbreak{}|\allowbreak{}\_\allowbreak{}X\allowbreak{}/\allowbreak{}(\allowbreak{}1\allowbreak{}-\allowbreak{}c\allowbreak{}h\allowbreak{}i\allowbreak{})\allowbreak{}.\allowbreak{}}
Retaining both terms in this mean decomposition and using (U28) yields
the additional full-row bound

\begin{verbatim}
||Chat(tau;x)||row
   <=[(1+255chi)/(1-chi)^2] exp[-3(1-chi)tau].        (U29a)
\end{verbatim}

The original bound(U29) remains available. Take the exact original
source radius and its derived constants

\begin{verbatim}
R1=45/4096,  chi(R1)=3/8,
C1=6184/25,  d1=15/8.
\end{verbatim}

Then (U29a) supplies

\begin{verbatim}
||Chat-C0-xi^2 C2-xi^4 C4||row
  <=(6184/25) exp(-15tau/8)
         (|xi|/R1)^6/[1-(|xi|/R1)^2], |xi|<R1.        (U29b)
\end{verbatim}

This second circle has its stated smaller radius and sharper prefactor
and decay. A scalar minimum of the two error bounds is permitted; no
minimum of noncommuting matrices is introduced. For the kth inverse
moment its prefactor is
{(\allowbreak{}6\allowbreak{}1\allowbreak{}8\allowbreak{}4\allowbreak{}/\allowbreak{}2\allowbreak{}5\allowbreak{})\allowbreak{}(\allowbreak{}8\allowbreak{}/\allowbreak{}1\allowbreak{}5\allowbreak{})\allowbreak{}\textasciicircum{}\allowbreak{}k\allowbreak{}.\allowbreak{}}
These are used in the strengthened numerical return M10.

\subsection{U8. Uniform analytic remainder and all inverse-energy
moments}\label{u8.-uniform-analytic-remainder-and-all-inverse-energy-moments}

Now restrict x\_p=zeta homogeneously. The original central link map
U\_(n,i) -\textgreater{}
{(\allowbreak{}-\allowbreak{}1\allowbreak{})\allowbreak{}\textasciicircum{}\allowbreak{}(\allowbreak{}s\allowbreak{}u\allowbreak{}m\allowbreak{}\_\allowbreak{}(\allowbreak{}j\allowbreak{}\textless{}\allowbreak{}i\allowbreak{})\allowbreak{}n\allowbreak{}\_\allowbreak{}j\allowbreak{})\allowbreak{}U\allowbreak{}\_\allowbreak{}(\allowbreak{}n\allowbreak{},\allowbreak{}i\allowbreak{})\allowbreak{}}
preserves Haar, K and every gauge constraint and sends all Wp to -Wp.
Both marked traces change sign, so Chat(tau;zeta) is even. Let
theta=\textbar xi\textbar/R0. Cauchy\textquotesingle s formula using
(U29), first at radius r\textless R0 and then letting r increase to R0,
proves in maximum absolute row-sum norm

\begin{verbatim}
||Chat(tau;xi)-C0(tau)-xi^2 C2(tau)-xi^4 C4(tau)||row
   <=512 exp(-3tau/2) theta^6/(1-theta^2),
   |xi|<R0, tau>=0.                                    (U30)
\end{verbatim}

The original coefficients C0,C2,C4 are exactly the inherited full heat
tables: uniqueness of the local analytic eigen/heat equation at x=0
gives the identical recurrence and ground-pole subtraction. Their source
and original time coordinates are unchanged.

For the original centered columns R and A=H-E0 put

\begin{verbatim}
G0=R*R,  Gk=kappa^k R*A^(-k)R (k>=1),
Phi=kappa A^(-1)R,  Phi*Phi=G2,  Phi*A Phi=kappa G1.
\end{verbatim}

The physical spectral theorem and (U30) justify all integrals, including
their complete tails:

\begin{verbatim}
Gk=int_0^infinity tau^(k-1) Chat(tau)d tau/(k-1)!,
||Gk-sum_(j=0)^2 xi^(2j) Gk_(2j)||row
   <=512(2/3)^k theta^6/(1-theta^2), k>=1.              (U31)
\end{verbatim}

For G0, (U30) at tau=0 supplies the same formula with k=0. A pole
a/(z+c)\^{}n contributes
{a\allowbreak{}(\allowbreak{}k\allowbreak{}+\allowbreak{}n\allowbreak{}-\allowbreak{}2\allowbreak{})\allowbreak{}!\allowbreak{}/\allowbreak{}[\allowbreak{}(\allowbreak{}k\allowbreak{}-\allowbreak{}1\allowbreak{})\allowbreak{}!\allowbreak{}(\allowbreak{}n\allowbreak{}-\allowbreak{}1\allowbreak{})\allowbreak{}!\allowbreak{}c\allowbreak{}\textasciicircum{}\allowbreak{}(\allowbreak{}k\allowbreak{}+\allowbreak{}n\allowbreak{}-\allowbreak{}1\allowbreak{})\allowbreak{}]\allowbreak{}.\allowbreak{}}
These maps use the original kappa in both the heat-time integral and
every inverse energy.

\subsection{U9. Spatial tails and the actual volume
limit}\label{u9.-spatial-tails-and-the-actual-volume-limit}

The preceding estimates also specify their spatial return. Set all
couplings outside a given finite list to zero. The scalar Hilbert tensor
factors on disjoint active link components factorize, as do the positive
vacuum and semigroup. Their restrictions to the original invariant
observables agree with that scalar calculation. A connected marked
correlation therefore has zero multivariate coefficient unless its
original marks and source faces lie in one edge-connected component. In
particular, if the original face-adjacency distance is d(p,q), a
degree-n coefficient connecting distinct marks needs
n\textgreater=d(p,q)-1. Both marks retain their original links.

At degree n, the same coefficient depends only on the n-neighborhood of
the union of marked links, where two links are adjacent when they share
an original plaquette. This follows either from the tensor factorization
or directly from (U8),(U17),(U25): every source insertion attaches its
original connected plaquette union through a differentiated edge, K and
its resolvent preserve link support, and all orders add to n. The
constant and disconnected products are retained until the connected
subtraction.

For two boxes containing the N-neighborhood of the same marks, every
homogeneous coefficient through N is identical.
Equation(U30)\textquotesingle s coefficient bound gives

\begin{verbatim}
|Chat_L-Chat_L'|
   <=1024 exp(-3tau/2) theta^(2(floor(N/2)+1))/(1-theta^2). (U32)
\end{verbatim}

The identical reasoning gives a row tail outside face-distance D by
omitting every coefficient with n\textless D-1; n is even. These are
explicit spatial bounds, with the physical link length a unaltered.

For clarity the limiting state and dynamics are also constructed. The
coefficient mean (U18) is analytic and bounded by
\textbar\textbar F\textbar\textbar X, and its Taylor coefficients depend
on finite original neighborhoods. Cauchy tails prove full-sequence
convergence mu\_L(F)-\textgreater mu(F) for every cylinder Fourier
polynomial on \textbar xi\textbar\textless R0. Real-coupling positivity
and mass one pass from the original vacua. For an arbitrary cylinder
function use the actual finite vertex Haar average P\_G on its incident
vertices; the equality mu\_L(F)=mu\_L(P\_G F) extends the limiting
physical functional consistently to all cylinder functions. Uniform
cylinder approximation then gives its gauge-invariant probability lift
on the product of original link groups. Both positivity and this gauge
extension are the original finite-volume ones.

Equation(U25) has norm at most3\textbar\textbar F\textbar\textbar X and
localized Taylor coefficients. Hence T\_L(tau)F converges uniformly,
first on cylinder Fourier polynomials and then on their continuous
uniform completion by the real Markov contraction. The resulting T is
positive, unital, strongly continuous, and preserves mu. Its semigroup
and symmetry laws pass by approximating the intermediate continuous
function uniformly with cylinder polynomials. Equations(U24)--(U26)
prove uniqueness of its invariant probability on the gauge-invariant
observable algebra: integrate the uniform mixing bound and use density.
This claim concerns physical observables; it does not assign a unique
gauge-variant probability from those observations alone.

The symmetric Markov semigroup has its self-adjoint generator A\_infty
on L2(mu)\_phys. The finite real bounds pass on cylinder functions by
the uniform convergence and then by L2 density, giving

\begin{verbatim}
A_infty|_(1 perp)>=kappa d(|xi|).                       (U33)
\end{verbatim}

On an original cylinder Fourier polynomial F, the functions (K-D\_v,L)F
converge in X to their infinite-label sum: at each source order they
involve finite neighborhoods, and the full drift tails are bounded by
the convergent local-source tail times
\textbar\textbar KF\textbar\textbar\_X. Passing to the semigroup
integral equation proves F belongs to the limiting generator domain with
A\_infty
{F\allowbreak{}=\allowbreak{}k\allowbreak{}a\allowbreak{}p\allowbreak{}p\allowbreak{}a\allowbreak{}(\allowbreak{}K\allowbreak{}-\allowbreak{}D\allowbreak{}\_\allowbreak{}v\allowbreak{},\allowbreak{}i\allowbreak{}n\allowbreak{}f\allowbreak{}t\allowbreak{}y\allowbreak{})\allowbreak{}F\allowbreak{}.\allowbreak{}}
In particular
{q\allowbreak{}\_\allowbreak{}i\allowbreak{}n\allowbreak{}f\allowbreak{}t\allowbreak{}y\allowbreak{}(\allowbreak{}F\allowbreak{},\allowbreak{}G\allowbreak{})\allowbreak{}=\allowbreak{}k\allowbreak{}a\allowbreak{}p\allowbreak{}p\allowbreak{}a\allowbreak{}}
mu Gamma(conjugate(F),G) on this cylinder core. This also fixes the
kinetic moment in the limiting original energy pairing.

All local time-ordered correlations pass through the same
product/semigroup limits. The row tails above construct bounded infinite
plaquette Gram operators on l2(P\_Z3), with the same bounds
(U29)--(U31). Their coefficientwise local limit and norm-bounded
convergence on finitely supported vectors give strong convergence on
that l2 space. This is a fixed-a spatial-volume construction on the
displayed coupling domain. It makes no assertion of an ultraviolet field
or finite continuum mass.

\subsection{U10. Primary sources and
scope}\label{u10.-primary-sources-and-scope}

G. Lumer and R. S. Phillips, Dissipative operators in a Banach space,
Pacific J. Math.11(1961),679--698, Theorem3.1, pp686--687: dense-domain
dissipativity plus actual range surjectivity gives the contraction
semigroup. The original theorem pages were inspected. The
coefficient-domain verification is (U22)--(U24).

P. Eymard, L\textquotesingle algebre de Fourier d\textquotesingle un
groupe localement compact, Bull. Soc. Math. France92(1964),181--236:
Fourier-algebra antecedent. The needed compact coefficient product and
domain estimates are proved here in (U3)--(U9); no claim is made to have
re-audited the complete historical paper.

D. Schuette, Zheng Weihong and C. J. Hamer, The Coupled Cluster Method
in Hamiltonian Lattice Field Theory, hep-lat/9603026:
exponential-vacuum, gauge-invariant character, and excitation
antecedents. The original constants and new heat row-sum return are
derived above. The predecessor heat catalogue and the point-of-use
source paths are in
{S\allowbreak{}O\allowbreak{}U\allowbreak{}R\allowbreak{}C\allowbreak{}E\allowbreak{}\_\allowbreak{}I\allowbreak{}N\allowbreak{}T\allowbreak{}A\allowbreak{}K\allowbreak{}E\allowbreak{}.\allowbreak{}j\allowbreak{}s\allowbreak{}o\allowbreak{}n\allowbreak{}.\allowbreak{}}
Numerical checks certify their declared finite algebra and endpoints,
not the written analytical arguments as a formal proof.

\section{Growing original plaquette families and their full physical
complement}\label{growing-original-plaquette-families-and-their-full-physical-complement}

21 September 2026. This is the quantitative return of
{V\allowbreak{}O\allowbreak{}L\allowbreak{}U\allowbreak{}M\allowbreak{}E\allowbreak{}\_\allowbreak{}U\allowbreak{}N\allowbreak{}I\allowbreak{}F\allowbreak{}O\allowbreak{}R\allowbreak{}M\allowbreak{}\_\allowbreak{}H\allowbreak{}E\allowbreak{}A\allowbreak{}T\allowbreak{}.\allowbreak{}m\allowbreak{}d\allowbreak{},\allowbreak{}}
with its source, graph, domain and semigroup proofs. Every finite open
box L\textgreater=2, spacing a\textgreater0, and coupling
g\^{}2\textgreater=16 is included below. All coefficient axes are
original plaquettes. Every Hilbert, energy and quotient pairing is
stated explicitly. The static and time-dependent coefficient tables are
the inherited original calculations, replayed unchanged; the
boundary-safe assembly and rational endpoints below are new
calculations. The analytic proof is written for review, not a formal
certificate.

\subsection{M1. Boundary-safe full-family
constants}\label{m1.-boundary-safe-full-family-constants}

Write R0=3/256, theta=xi/R0, and
{t\allowbreak{}6\allowbreak{}=\allowbreak{}t\allowbreak{}h\allowbreak{}e\allowbreak{}t\allowbreak{}a\allowbreak{}\textasciicircum{}\allowbreak{}6\allowbreak{}/\allowbreak{}(\allowbreak{}1\allowbreak{}-\allowbreak{}t\allowbreak{}h\allowbreak{}e\allowbreak{}t\allowbreak{}a\allowbreak{}\textasciicircum{}\allowbreak{}2\allowbreak{})\allowbreak{}.\allowbreak{}}
The complete marked heat tables have 199 connected multisets containing
the anchor face (0,0,0;0,1): one degree-zero self term; one self and
twelve adjacent-pair terms at degree two; and one self, twenty-four
ordered repeated-pair, 138 path, twelve common-edge, eight corner and
two cube terms at degree four. Their 559 marked-row contributions are
all retained.

For each contribution, integrate its actual partial fractions as in U31,
or take the original time-zero value. Sum absolute values before
different original supports can cancel. Every finite box contains a
subset of these original contributions at a specified anchor. Thus the
following are bounds for every original row, including every boundary
row, independent of box size:

{\def\LTcaptype{none} % do not increment counter
\begin{longtable}[]{@{}lrr@{}}
\toprule\noalign{}
Original matrix & degree2 bound & degree4 bound \\
\midrule\noalign{}
\endhead
\bottomrule\noalign{}
\endlastfoot
G0 & 149/468 &
{2\allowbreak{}3\allowbreak{}6\allowbreak{}1\allowbreak{}9\allowbreak{}9\allowbreak{}4\allowbreak{}0\allowbreak{}7\allowbreak{}3\allowbreak{}/\allowbreak{}4\allowbreak{}6\allowbreak{}9\allowbreak{}1\allowbreak{}4\allowbreak{}9\allowbreak{}4\allowbreak{}0\allowbreak{}8\allowbreak{}0\allowbreak{}} \\
G1 & 97/468 &
{3\allowbreak{}7\allowbreak{}6\allowbreak{}8\allowbreak{}8\allowbreak{}2\allowbreak{}6\allowbreak{}9\allowbreak{}1\allowbreak{}/\allowbreak{}9\allowbreak{}3\allowbreak{}8\allowbreak{}2\allowbreak{}9\allowbreak{}8\allowbreak{}8\allowbreak{}1\allowbreak{}6\allowbreak{}} \\
G2 & 73313/657072 &
{9\allowbreak{}8\allowbreak{}2\allowbreak{}0\allowbreak{}6\allowbreak{}9\allowbreak{}7\allowbreak{}1\allowbreak{}8\allowbreak{}9\allowbreak{}6\allowbreak{}3\allowbreak{}8\allowbreak{}3\allowbreak{}3\allowbreak{}/\allowbreak{}3\allowbreak{}9\allowbreak{}1\allowbreak{}9\allowbreak{}1\allowbreak{}8\allowbreak{}0\allowbreak{}3\allowbreak{}2\allowbreak{}4\allowbreak{}5\allowbreak{}5\allowbreak{}0\allowbreak{}4\allowbreak{}0\allowbreak{}0\allowbreak{}} \\
Kobs & 187/468 &
{5\allowbreak{}8\allowbreak{}6\allowbreak{}6\allowbreak{}6\allowbreak{}8\allowbreak{}4\allowbreak{}2\allowbreak{}1\allowbreak{}/\allowbreak{}1\allowbreak{}5\allowbreak{}6\allowbreak{}3\allowbreak{}8\allowbreak{}3\allowbreak{}1\allowbreak{}3\allowbreak{}6\allowbreak{}0\allowbreak{}} \\
J3=G0-6G1+9G2 & 6779/73008 &
{1\allowbreak{}5\allowbreak{}2\allowbreak{}3\allowbreak{}3\allowbreak{}8\allowbreak{}6\allowbreak{}7\allowbreak{}4\allowbreak{}0\allowbreak{}0\allowbreak{}5\allowbreak{}9\allowbreak{}8\allowbreak{}9\allowbreak{}/\allowbreak{}4\allowbreak{}3\allowbreak{}5\allowbreak{}4\allowbreak{}6\allowbreak{}4\allowbreak{}4\allowbreak{}8\allowbreak{}0\allowbreak{}5\allowbreak{}0\allowbreak{}5\allowbreak{}6\allowbreak{}0\allowbreak{}0\allowbreak{}} \\
\end{longtable}
}

Call these numbers B\_(k,2),B\_(k,4), and similarly B\_K,B\_J. The exact
transport for every contribution is the original signed coordinate
permutation and translation. The independent audit constructs the 199
original multisets without calling that transport producer, verifies its
inverse coordinate formula, and reintegrates each original pole
directly. The exact row entries of the existing full L2 matrix lie below
these bounds; no assertion that the bulk signed row is the largest
boundary row is used.

The kinetic matrix is
{K\allowbreak{}o\allowbreak{}b\allowbreak{}s\allowbreak{}=\allowbreak{}k\allowbreak{}a\allowbreak{}p\allowbreak{}p\allowbreak{}a\allowbreak{}\textasciicircum{}\allowbreak{}(\allowbreak{}-\allowbreak{}1\allowbreak{})\allowbreak{}R\allowbreak{}*\allowbreak{}A\allowbreak{}R\allowbreak{}=\allowbreak{}\textless{}\allowbreak{}G\allowbreak{}a\allowbreak{}m\allowbreak{}m\allowbreak{}a\allowbreak{}(\allowbreak{}W\allowbreak{}p\allowbreak{},\allowbreak{}W\allowbreak{}q\allowbreak{})\allowbreak{}\textgreater{}\allowbreak{}r\allowbreak{}h\allowbreak{}o\allowbreak{}.\allowbreak{}}
Only p=q or faces sharing an original edge can contribute. At p=q the
exact function is
{4\allowbreak{}-\allowbreak{}W\allowbreak{}p\allowbreak{}\textasciicircum{}\allowbreak{}2\allowbreak{}=\allowbreak{}3\allowbreak{}-\allowbreak{}c\allowbreak{}h\allowbreak{}i\allowbreak{}\_\allowbreak{}1\allowbreak{}(\allowbreak{}W\allowbreak{}p\allowbreak{})\allowbreak{}.\allowbreak{}}
Its trace-coefficient norm is at most30: the original four-edge spin-one
contraction has norm27, as can be obtained by contracting its
three-dimensional indices, or bounded by the original tensor coefficient
with one turn. For adjacent faces the direct derivative estimate U4
gives norm at most3*(8/2)\^{}2=48. There are at most twelve neighbors.
Therefore, on the full complex source disk,

\begin{verbatim}
||Kobs(zeta)||row<=30+12*48=606.                        (M1)
\end{verbatim}

The spin-one contraction used for30 has the same four-link index formula
as U10 with the SU(2) invariant dual matrix in dimension3; its nonzero
singular values are3 with total trace norm27. In the original
magnetic-weight basis its fixed dual matrix is C\_(mn)=(-1)\^{}(j-m)
delta\_(m,-n), m,n=-j,...,j. Thus pi\_j(U\^{}(-1))=C pi\_j(U)\^{}T
C\^{}(-1); both coefficient-side unitary factors are retained. For the
trace of four general d-by-d factors with its last two transposed, the
coefficient has A\_(b,c,c,e ; a,b,e,a) +=1. Direct multiplication gives
(A* A)\^{}2=d\^{}2 A* A and Tr(A* A)=d\^{}4, so its trace norm is
d\^{}3. Taking d=3 proves27. Its inverse-link signs are represented by
the stated dual matrix, not omitted. The accompanying finite check
verifies the dimension2 and dimension3 contraction directly.
Alternatively the bound4+8\^{}2=68 would give644 and slightly wider
kinetic endpoints.

Combining U30--U31, M1, and the table proves

\begin{verbatim}
||Gk-3^(-k)I|| <= xi^2 B_(k,2)+xi^4 B_(k,4)
                     +512(2/3)^k t6 = w_k, k=0,1,2,
||Kobs-3I|| <= xi^2 B_(K,2)+xi^4 B_(K,4)+606 t6=w_K,
0<=J3<=beta I,
beta=xi^2 B_(J,2)+xi^4 B_(J,4)+4608 t6.                 (M2)
\end{verbatim}

The whole-matrix bound follows from equality of row and column absolute
norm for the original symmetric real matrices, followed by the l2 Schur
bound. The
{r\allowbreak{}e\allowbreak{}m\allowbreak{}a\allowbreak{}i\allowbreak{}n\allowbreak{}d\allowbreak{}e\allowbreak{}r\allowbreak{}4\allowbreak{}6\allowbreak{}0\allowbreak{}8\allowbreak{}=\allowbreak{}5\allowbreak{}1\allowbreak{}2\allowbreak{}+\allowbreak{}6\allowbreak{}*\allowbreak{}5\allowbreak{}1\allowbreak{}2\allowbreak{}*\allowbreak{}(\allowbreak{}2\allowbreak{}/\allowbreak{}3\allowbreak{})\allowbreak{}+\allowbreak{}9\allowbreak{}*\allowbreak{}5\allowbreak{}1\allowbreak{}2\allowbreak{}*\allowbreak{}(\allowbreak{}2\allowbreak{}/\allowbreak{}3\allowbreak{})\allowbreak{}\textasciicircum{}\allowbreak{}2\allowbreak{}}
includes all three original moment errors. The matrix J3 is positive
because it is exactly (R-3Phi)*(R-3Phi), with Phi=kappa A\^{}(-1)R.

\subsection{M2. Actual numerical bounds for every original
box}\label{m2.-actual-numerical-bounds-for-every-original-box}

For xi\textless=1/1024, equivalently g\^{}2\textgreater=16,
theta\textless=1/12. Every term in (M2) is increasing in xi. Evaluation
at the endpoint by exact rational arithmetic gives

\begin{verbatim}
w_0<173/10^6,     w_1<116/10^6,
w_2<77/10^6,      w_K<205/10^6,
beta<1555/10^6.                                        (M3)
\end{verbatim}

The sharper fractions, not these rounded values, are retained in
{g\allowbreak{}e\allowbreak{}n\allowbreak{}e\allowbreak{}r\allowbreak{}a\allowbreak{}t\allowbreak{}e\allowbreak{}d\allowbreak{}/\allowbreak{}u\allowbreak{}n\allowbreak{}i\allowbreak{}f\allowbreak{}o\allowbreak{}r\allowbreak{}m\allowbreak{}\_\allowbreak{}c\allowbreak{}o\allowbreak{}n\allowbreak{}s\allowbreak{}t\allowbreak{}a\allowbreak{}n\allowbreak{}t\allowbreak{}s\allowbreak{}.\allowbreak{}j\allowbreak{}s\allowbreak{}o\allowbreak{}n\allowbreak{}.\allowbreak{}}
Set

\begin{verbatim}
l0=1-w0, u0=1+w0,
l1=1/3-w1, u1=1/3+w1,
l2=1/9-w2, u2=1/9+w2,
lK=3-wK, uK=3+wK,                                     (M4)
\end{verbatim}

using the endpoint rational w values. They are positive. The full
physical source proof U26 gives A\textgreater=kappa d I with

\begin{verbatim}
d=117/40,                                             (M5)
\end{verbatim}

since at xi=1/1024 its square-root argument
{i\allowbreak{}s\allowbreak{}1\allowbreak{}1\allowbreak{}/\allowbreak{}1\allowbreak{}2\allowbreak{}\textgreater{}\allowbreak{}(\allowbreak{}1\allowbreak{}9\allowbreak{}/\allowbreak{}2\allowbreak{}0\allowbreak{})\allowbreak{}\textasciicircum{}\allowbreak{}2\allowbreak{}.\allowbreak{}}
This is a lower bound on the entire centered physical space, proved
without selecting an observation family. All physical dimensions and
original energies in (M3)--(M5) remain.

\subsection{M3. The original observation and its inverse-energy
family}\label{m3.-the-original-observation-and-its-inverse-energy-family}

Let R:C\^{}M-\textgreater H0 have columns
{r\allowbreak{}p\allowbreak{}=\allowbreak{}(\allowbreak{}W\allowbreak{}p\allowbreak{}-\allowbreak{}\textless{}\allowbreak{}W\allowbreak{}p\allowbreak{}\textgreater{}\allowbreak{}r\allowbreak{}h\allowbreak{}o\allowbreak{})\allowbreak{}p\allowbreak{}s\allowbreak{}i\allowbreak{}.\allowbreak{}}
On its literal coefficient coordinates put

\begin{verbatim}
Phi=kappa A^(-1)R,
G0=R*R, G1=R*Phi, G2=Phi*Phi, E=Phi*A Phi=kappa G1.
\end{verbatim}

The inequalities above prove R and Phi injective and their ranges
closed. The original observation O=R* has exact left inverse on the
primitive family

\begin{verbatim}
W=G1^(-1)R*,  W Phi=I.                                (M6)
\end{verbatim}

The state projection P\_R=R G0\^{}(-1)R* retains

\begin{verbatim}
||P_R Phi c||^2=c*G1 G0^(-1)G1 c.
\end{verbatim}

Using l1,l2,u0,u2 from the original matrices proves

\begin{verbatim}
||P_R Phi c||^2 >= [l1^2/(u0 u2)]||Phi c||^2
                  >(624/625)||Phi c||^2.              (M7)
\end{verbatim}

This is the inverse-power
{o\allowbreak{}b\allowbreak{}s\allowbreak{}e\allowbreak{}r\allowbreak{}v\allowbreak{}a\allowbreak{}t\allowbreak{}i\allowbreak{}o\allowbreak{}n\allowbreak{}/\allowbreak{}m\allowbreak{}i\allowbreak{}n\allowbreak{}i\allowbreak{}m\allowbreak{}u\allowbreak{}m\allowbreak{}-\allowbreak{}s\allowbreak{}e\allowbreak{}c\allowbreak{}t\allowbreak{}i\allowbreak{}o\allowbreak{}n\allowbreak{}}
mechanism of the RH workbench\textquotesingle s IK9--12, instantiated
here by the exact matrix G1\^{}(-1)R* and the calculated physical Grams.
No arithmetic constant or conductor coefficient enters (M7).

The correction between the minimum-state and minimum-energy sections of
this observation is h\_R=(I-P\_R)Phi. Its original energy is

\begin{verbatim}
h_R*A h_R
   =kappa[G1 G0^(-1)Kobs G0^(-1)G1-G1].               (M8)
\end{verbatim}

Indeed expansion has two cross terms -kappa G1 and the positive section
term; both are retained. Also q\_A(Phi c,h\_R d)=kappa\textless Rc,h\_R
d\textgreater=0. Consequently

\begin{verbatim}
0<=h_R*A h_R
   <=[u1 uK/l0^2-1] E <(1/1250)E.                    (M9)
\end{verbatim}

\subsection{M4. The complete surrounding physical
space}\label{m4.-the-complete-surrounding-physical-space}

The state projection onto the primitive family is instead

\begin{verbatim}
P=Phi G2^(-1)Phi*, Q=I-P,
B=QR:C^M->QH0.                                        (M10)
\end{verbatim}

It gives the exact leakage Gram

\begin{verbatim}
B*B=G0-G1 G2^(-1)G1.                                  (M11)
\end{verbatim}

For each c, P Rc is the state minimum over Phi x. Inserting the
particular original-coordinate trial x=3c gives

\begin{verbatim}
0<=B*B<=(R-3Phi)*(R-3Phi)=J3<=beta I.                  (M12)
\end{verbatim}

No commutativity of G0,G1,G2 is used in this argument.

Every original form vector decomposes as Phi c+h, h in QH0. Its state
norm and complete energy are

\begin{verbatim}
||Phi c+h||^2=c*G2 c+||h||^2,
q_A(Phi c+h)=c*E c+2kappa Re<Bc,h>+q_A(h).             (M13)
\end{verbatim}

Since q\_A(h)\textgreater=kappa
d\textbar\textbar h\textbar\textbar\^{}2, the complete mixed term
satisfies

\begin{verbatim}
|2kappa Re<Bc,h>|
   <=2 sqrt(beta/(d l1)) sqrt(c*E c . q_A(h)).          (M14)
\end{verbatim}

The actual endpoint fractions give beta/(d l1)\textless1/625. Thus every
vector in the entire physical form domain satisfies the coupled
comparison

\begin{verbatim}
(24/25)[c*E c+q_A(h)]
   <=q_A(Phi c+h)
   <=(26/25)[c*E c+q_A(h)].                            (M15)
\end{verbatim}

The off-diagonal term is controlled, not set equal to zero. This is a
full-space comparison, not a Ritz lower bound on a selected finite
subspace.

\subsection{M5. The original complementary inverse and complete Schur
loss}\label{m5.-the-original-complementary-inverse-and-complete-schur-loss}

A Phi=kappa R is bounded as a coefficient map. Hence AP is bounded and
PA is its adjoint. Subtract the bounded self-adjoint off-diagonal
QAP+PAQ from A. The resulting self-adjoint operator has reducing
subspaces PH0,QH0. Its Q restriction is

\begin{verbatim}
D_Q=QAQ, Dom(D_Q)=Dom(A) intersect QH0,
D_Q>=kappa d I.                                       (M16)
\end{verbatim}

The inverse is therefore supplied by the full physical estimate,
including every vector of the complementary space. Minimizing M13 over
that whole space gives

\begin{verbatim}
h_min(c)=-kappa D_Q^(-1)Bc,
E_full=E-kappa^2 B*D_Q^(-1)B.                          (M17)
\end{verbatim}

Testing the residual against all complementary form vectors proves this
is the actual minimum, not only a finite-column stationary point. The
full loss and restored state metric obey

\begin{verbatim}
0<=kappa^2 B*D_Q^(-1)B
   <=(kappa beta/d)I <(1/625)E,
(624/625)E<E_full<=E,                                 (M18)

G_restored=G2+kappa^2 B*D_Q^(-2)B,
G2<=G_restored<(601/600)G2.                            (M19)
\end{verbatim}

The second bound uses beta/(d\^{}2 l2)\textless1/600. All signs, kappa
factors, original Grams and complementary feedback are retained. The
graph map c-\textgreater Phi c+h\_min(c) has these exact energy and
state forms. These estimates address the full complementary coupling
left by the predecessor heat construction, with constants independent of
both the number of plaquettes and the exterior box.

\subsection{M6. Fixed supports, their quotients, and a common physical
heat
horizon}\label{m6.-fixed-supports-their-quotients-and-a-common-physical-heat-horizon}

For original face sets F subset G put V\_F=Phi(C\^{}F), U\_F=A
V\_F=kappa R(C\^{}F). The original complexes are

\begin{verbatim}
V_F --A--> H0 --0-->0.
\end{verbatim}

Their support transitions are inclusion in degree0 and identity in
degree1. The exact transported kernel is

\begin{verbatim}
V_G/V_F -> ker[H0/U_F -> H0/U_G], [h]->[Ah].            (M20)
\end{verbatim}

Its inverse sends {[}Ah{]} to {[}h{]}. Injectivity of A on H0 proves the
inverse is independent of representatives. Its coefficient-energy
minimum, with J=G\textbackslash F, is

\begin{verbatim}
c_F=-E_FF^(-1)E_FJ c_J,
Q_E=E_JJ-E_JF E_FF^(-1)E_FJ.                           (M21)
\end{verbatim}

The removed old-support primitive Phi\_F c\_F and both mixed entries
remain. The bounds kappa l1 I\textless=E\textless=kappa u1 I imply
exactly the same bounds for every such quotient minimum in its original
J coordinates. They follow by minimizing the two complete inequalities
on the identical affine fiber.

For T\textgreater0 the actual heat columns and forcing defect are

\begin{verbatim}
Phi_T=(I-exp(-TA))Phi=kappa int_0^T exp(-tA)Rdt,
A Phi_T=kappa R-kappa exp(-TA)R.                       (M22)
\end{verbatim}

The omitted primitive is -exp(-TA)Phi. On the complete centered physical
space, I-exp(-TA) has inverse
{s\allowbreak{}u\allowbreak{}m\allowbreak{}\_\allowbreak{}(\allowbreak{}n\allowbreak{}\textgreater{}\allowbreak{}=\allowbreak{}0\allowbreak{})\allowbreak{}e\allowbreak{}x\allowbreak{}p\allowbreak{}(\allowbreak{}-\allowbreak{}n\allowbreak{}T\allowbreak{}A\allowbreak{})\allowbreak{},\allowbreak{}}
also on the graph domain, and commutes with A. Thus it gives a cochain
isomorphism of every original support diagram. Applying the Split Zero
support reconstruction retains its support label and receiving zero
together with this actual primitive.

Set epsilon=exp(-kappa d T). Expanding all three heat Gram terms and
then applying spectral calculus yields

\begin{verbatim}
(1-epsilon)^2 G2<=Phi_T*Phi_T<=G2,
(1-epsilon)^2 E<=Phi_T*A Phi_T<=E.                     (M23)
\end{verbatim}

Taking minima over the same fixed fibers carries these inequalities
through every finite sequence of original restrictions and quotients,
without multiplying a constant for each step. This is precisely the
HM19--24 minimum-fiber argument of the RH source, whose coefficient maps
remain fixed before the metric comparison. Each rank remains in the
determinant bound:

\begin{verbatim}
|L_T-L|<=2 B[-log(1-epsilon)]<=2B epsilon/(1-epsilon),
B=sum_j |a_j| rank_j, L=sum_j a_j log det Q_j.          (M24)
\end{verbatim}

Rank-zero terms keep determinant1. For four signed returns of rank at
most240, the single physical horizon T=10/kappa=5a/g\^{}2 gives an error
below4*10\^{}(-10), in every exterior volume. For ranks that grow with
M, a specified tolerance eta\textgreater0 is reached at the explicit
original physical time

\begin{verbatim}
T=(kappa d)^(-1)log(1+2B/eta).                         (M25)
\end{verbatim}

Thus no growing determinant rank has been discarded. The original
kappa\^{}rank factors cancel only between the two determinants in the
identical energy frame.

\subsection{M7. The actual opposite-face result in every exterior
box}\label{m7.-the-actual-opposite-face-result-in-every-exterior-box}

Keep p=(0,0,0;0,1), q=(0,0,1;0,1), xi=10\^{}(-10), g\^{}2=50000 and
kappa=100000/a. Their inherited complete coefficient h4(tau) has zero
degrees0 and2. All original supports contributing to degree4 lie inside
the L=2 box; hence the same h4 applies in every L\textgreater=2. The new
uniform tail U30 gives

\begin{verbatim}
85910/10^7 < C_pq(1/kappa;xi)/xi^4 < 85920/10^7.        (M26)
\end{verbatim}

The previously proved exact coefficient lower bound
h4(tau)\textgreater=5 exp(-3tau)/648, tau\textgreater=1, is kept. On
1\textless=tau\textless=3 the ratio of U30\textquotesingle s full
remainder to that lower bound is at most

\begin{verbatim}
(512*648/5) xi^2 R0^(-6)
   exp(9/2)/[1-(xi/R0)^2] <24/1000.                   (M27)
\end{verbatim}

Thus the actual connected heat correlation is strictly positive
throughout
{1\allowbreak{}/\allowbreak{}k\allowbreak{}a\allowbreak{}p\allowbreak{}p\allowbreak{}a\allowbreak{}\textless{}\allowbreak{}=\allowbreak{}t\allowbreak{}\textless{}\allowbreak{}=\allowbreak{}3\allowbreak{}/\allowbreak{}k\allowbreak{}a\allowbreak{}p\allowbreak{}p\allowbreak{}a\allowbreak{},\allowbreak{}}
uniformly over the original exterior volume. The result also passes to
the constructed spatial-volume limit U33. These statements use all
higher orders, not only the sign of h4.

\subsection{M8. Infinite plaquette families and the original continuum
path}\label{m8.-infinite-plaquette-families-and-the-original-continuum-path}

The row-tail and local-limit constructions U9 give bounded operators
G0,G1,G2,Kobs on l2 of all original plaquettes, with exactly the same
numerical intervals (M2)--(M5). On finitely supported c, R c is the
original centered local observable in L2(mu)\_phys. In this
infinite-family statement Kobs denotes the closed-form pairing
{k\allowbreak{}a\allowbreak{}p\allowbreak{}p\allowbreak{}a\allowbreak{}\textasciicircum{}\allowbreak{}(\allowbreak{}-\allowbreak{}1\allowbreak{})\allowbreak{}q\allowbreak{}\_\allowbreak{}(\allowbreak{}A\allowbreak{}\_\allowbreak{}i\allowbreak{}n\allowbreak{}f\allowbreak{}t\allowbreak{}y\allowbreak{})\allowbreak{}(\allowbreak{}R\allowbreak{}.\allowbreak{},\allowbreak{}R\allowbreak{}.\allowbreak{})\allowbreak{},\allowbreak{}}
whose boundedness follows from U33 and the cylinder-generator identity
after it. It does not require asserting that A\_infty R is a bounded
operator on all l2 coefficients. G0 proves this map extends boundedly
and is bounded below. Define Phi=kappa A\_infty\^{}(-1)R. Then G1,G2 are
its original mixed and state Grams by the convergent physical heat
integrals. They are bounded and bounded below, so P=Phi G2\^{}(-1)Phi*
is again an actual orthogonal projection.

A\_infty Phi=kappa R is bounded; hence A\_infty P and PA\_infty have the
same bounded off-diagonal argument as M16. All maps, complementary
minima and inequalities M6--M19 therefore pass to the complete infinite
plaquette coefficient space and its full physical complement. This
passage uses proved bounded operators, not a formal limit of inverses
without a lower bound.

At fixed admissible g,a this controls the spatial-volume limit. For the
original simultaneous path a\_n=a0*2\^{}(-n),
{g\allowbreak{}\_\allowbreak{}n\allowbreak{}\textasciicircum{}\allowbreak{}2\allowbreak{}=\allowbreak{}1\allowbreak{}/\allowbreak{}(\allowbreak{}g\allowbreak{}0\allowbreak{}\textasciicircum{}\allowbreak{}(\allowbreak{}-\allowbreak{}2\allowbreak{})\allowbreak{}+\allowbreak{}b\allowbreak{}e\allowbreak{}t\allowbreak{}a\allowbreak{}}
n log2), the source disk used in U13 is entered exactly when
c\_n=g0\^{}(-2)+beta n log2\textless sqrt(3)/8. The narrower numerical
metric domain is c\_n\textless=1/16. For beta\textgreater0 the path
eventually leaves each of them. The current result removes the 1/M heat
radius and the unbounded observation-rank loss on a specified
strong-coupling domain; it does not establish a smooth four-dimensional
continuum field or finite continuum mass. No universal conclusion is
inferred merely from a support kernel or an auxiliary norm.

\subsection{M9. What is new in this
continuation}\label{m9.-what-is-new-in-this-continuation}

The older elementary local source argument U3--U6 has been rederived at
its own conservative radius. The new calculation combines its full
scalar return with the Poisson identity U21 to bound an entire
heat-matrix row, not just one entry. This produces the original
box-independent analytic remainder U30--U31. The 199-support absolute
assembly then gives original
{s\allowbreak{}t\allowbreak{}a\allowbreak{}t\allowbreak{}e\allowbreak{}/\allowbreak{}r\allowbreak{}e\allowbreak{}s\allowbreak{}p\allowbreak{}o\allowbreak{}n\allowbreak{}s\allowbreak{}e\allowbreak{}/\allowbreak{}k\allowbreak{}i\allowbreak{}n\allowbreak{}e\allowbreak{}t\allowbreak{}i\allowbreak{}c\allowbreak{}}
metrics for growing families, and M12--M19 controls their full
complementary coupling. The spatial and physical maps are explicit
throughout. No larger coupling threshold or previously unaudited
historical theorem is assigned a new verification status.

\subsection{M10. Sharper original-Haar return and broader family
domains}\label{m10.-sharper-original-haar-return-and-broader-family-domains}

The additional bound U29b gives the same complete inverse-moment errors
with the smaller circle R1=45/4096 and prefactors
{(\allowbreak{}6\allowbreak{}1\allowbreak{}8\allowbreak{}4\allowbreak{}/\allowbreak{}2\allowbreak{}5\allowbreak{})\allowbreak{}(\allowbreak{}8\allowbreak{}/\allowbreak{}1\allowbreak{}5\allowbreak{})\allowbreak{}\textasciicircum{}\allowbreak{}k\allowbreak{}.\allowbreak{}}
For a fixed real x in that disk put
{t\allowbreak{}0\allowbreak{}=\allowbreak{}(\allowbreak{}x\allowbreak{}/\allowbreak{}R\allowbreak{}0\allowbreak{})\allowbreak{}\textasciicircum{}\allowbreak{}6\allowbreak{}/\allowbreak{}[\allowbreak{}1\allowbreak{}-\allowbreak{}(\allowbreak{}x\allowbreak{}/\allowbreak{}R\allowbreak{}0\allowbreak{})\allowbreak{}\textasciicircum{}\allowbreak{}2\allowbreak{}]\allowbreak{},\allowbreak{}}
{t\allowbreak{}1\allowbreak{}=\allowbreak{}(\allowbreak{}x\allowbreak{}/\allowbreak{}R\allowbreak{}1\allowbreak{})\allowbreak{}\textasciicircum{}\allowbreak{}6\allowbreak{}/\allowbreak{}[\allowbreak{}1\allowbreak{}-\allowbreak{}(\allowbreak{}x\allowbreak{}/\allowbreak{}R\allowbreak{}1\allowbreak{})\allowbreak{}\textasciicircum{}\allowbreak{}2\allowbreak{}]\allowbreak{}.\allowbreak{}}
Define the actual scalar error bounds

\begin{verbatim}
w_k^sharp=x^2 B_(k,2)+x^4 B_(k,4)
   +min(512(2/3)^k t0,(6184/25)(8/15)^k t1), k=0,1,2,
beta^sharp=x^2 B_(J,2)+x^4 B_(J,4)
   +min(4608t0,(6184/25)[1+6(8/15)+9(8/15)^2]t1).
\end{verbatim}

The kinetic w\_K from M2 is retained. Each scalar expression is
increasing in x. Replacing the endpoint l and u values in the already
proved M7--M19 inequalities by these evaluated sharper endpoints gives
the following complete original-family bounds, all uniform over
L\textgreater=2 and a\textgreater0:

{\def\LTcaptype{none} % do not increment counter
\begin{longtable}[]{@{}
  >{\raggedright\arraybackslash}p{(\linewidth - 8\tabcolsep) * \real{0.1579}}
  >{\raggedleft\arraybackslash}p{(\linewidth - 8\tabcolsep) * \real{0.2105}}
  >{\raggedleft\arraybackslash}p{(\linewidth - 8\tabcolsep) * \real{0.2105}}
  >{\raggedleft\arraybackslash}p{(\linewidth - 8\tabcolsep) * \real{0.2105}}
  >{\raggedleft\arraybackslash}p{(\linewidth - 8\tabcolsep) * \real{0.2105}}@{}}
\toprule\noalign{}
\begin{minipage}[b]{\linewidth}\raggedright
g\^{}2 at least
\end{minipage} & \begin{minipage}[b]{\linewidth}\raggedleft
observed fraction in M7 exceeds
\end{minipage} & \begin{minipage}[b]{\linewidth}\raggedleft
E\_full/E exceeds
\end{minipage} & \begin{minipage}[b]{\linewidth}\raggedleft
G\_restored/G2 is less than
\end{minipage} & \begin{minipage}[b]{\linewidth}\raggedleft
h\_R energy / E is less than
\end{minipage} \\
\midrule\noalign{}
\endhead
\bottomrule\noalign{}
\endlastfoot
10 & 39/50 & 18/25 & 33/25 & 1/6 \\
12 & 97/100 & 243/250 & 103/100 & 1/60 \\
16 & 999/1000 & 999/1000 & 1001/1000 & 1/1900 \\
\end{longtable}
}

Here every entry with matrix forms means the displayed Loewner
inequality in the same original coefficient coordinates. The full-space
gap lower constants used in those three rows are respectively
{d\allowbreak{}=\allowbreak{}1\allowbreak{}4\allowbreak{}/\allowbreak{}5\allowbreak{},\allowbreak{}7\allowbreak{}2\allowbreak{}/\allowbreak{}2\allowbreak{}5\allowbreak{},\allowbreak{}1\allowbreak{}1\allowbreak{}7\allowbreak{}/\allowbreak{}4\allowbreak{}0\allowbreak{}}
times kappa. Each follows directly from U26 by squaring its positive
rational comparison. The file
{g\allowbreak{}e\allowbreak{}n\allowbreak{}e\allowbreak{}r\allowbreak{}a\allowbreak{}t\allowbreak{}e\allowbreak{}d\allowbreak{}/\allowbreak{}s\allowbreak{}h\allowbreak{}a\allowbreak{}r\allowbreak{}p\allowbreak{}\_\allowbreak{}c\allowbreak{}o\allowbreak{}n\allowbreak{}s\allowbreak{}t\allowbreak{}a\allowbreak{}n\allowbreak{}t\allowbreak{}s\allowbreak{}.\allowbreak{}j\allowbreak{}s\allowbreak{}o\allowbreak{}n\allowbreak{}}
contains every rational input, interval endpoint and returned fraction;
the independent auditor recomputes each from the original marked
coefficients.

At g\^{}2\textgreater=16 the absolute Gram widths obey

\begin{verbatim}
||G0-I||<124/10^6,
||G1-I/3||<66/10^6,
||G2-I/9||<36/10^6,
||Kobs-3I||<205/10^6,
B*B< (832/10^6) I.
\end{verbatim}

The exact unrounded beta\^{}sharp/(d l1\^{}sharp) is smaller than
1/34\^{}2. Accordingly the full mixed-space comparison M15 sharpens to

\begin{verbatim}
(33/34)[c*E c+q_A(h)] <= q_A(Phi c+h)
                        <=(35/34)[c*E c+q_A(h)].       (M28)
\end{verbatim}

The full complementary minimization loses less than 1/1000 of E and adds
less than 1/1000 of G2 to the restored state norm. No complementary
vector or source-kernel direction is removed. The same bounded-operator
argument in M8 carries these stronger constants to the infinite original
plaquette family. The earlier M2--M27 endpoints remain valid and have
their own exact records; this paragraph retains an additional,
explicitly proved original-Haar estimate rather than silently replacing
their coefficients.
\end{document}
